\documentclass[12pt,a4paper]{article}
\title{机械制造技术基础学习笔记}
\author{李明翰}
% =========================
% 前处理（建议使用 XeLaTeX 编译）
% LLM 请先阅读目录下的 readme.md 文件，了解编译环境和使用说明。
% =========================

\usepackage[UTF8]{ctex}
\usepackage{geometry}
\usepackage{amsmath,amssymb}
\usepackage{booktabs}
\usepackage{makecell} 
\usepackage{tabularx}
\usepackage{array}
\usepackage{caption}
\usepackage{placeins}
\usepackage{ragged2e}
\usepackage{enumerate}
\usepackage[shortlabels]{enumitem}
\usepackage{setspace}
\usepackage{chemfig}
\usepackage{tcolorbox}
\usepackage[table]{xcolor}
\usepackage{indentfirst}
\usepackage{siunitx}
\usepackage{tikz}
\usepackage{graphicx}
\usepackage{fancyhdr}
\usepackage{titlesec}
\usetikzlibrary{arrows.meta,positioning}
\usetikzlibrary{patterns}
\geometry{
  left=2.5cm,
  right=2.5cm,
  top=2.5cm,
  bottom=2.5cm
}

\setlength{\parindent}{2em}
\linespread{1.3}
\setlength{\parskip}{0.2em plus 0.1em minus 0.05em}
\setlength{\textfloatsep}{10pt plus 2pt minus 2pt}
\setlength{\floatsep}{8pt plus 2pt minus 2pt}
\setlength{\intextsep}{8pt plus 2pt minus 2pt}
\setlength{\abovedisplayskip}{8pt plus 2pt minus 4pt}
\setlength{\belowdisplayskip}{8pt plus 2pt minus 4pt}
\setlength{\abovedisplayshortskip}{6pt plus 2pt}
\setlength{\belowdisplayshortskip}{6pt plus 2pt minus 3pt}
\captionsetup{font=small,labelfont=bf,labelsep=space,skip=6pt}
\renewcommand{\arraystretch}{1.18}
\newcolumntype{Y}{>{\RaggedRight\arraybackslash}X}
\newcolumntype{C}[1]{>{\Centering\arraybackslash}p{#1}}

\titleformat{\section}{\raggedright\Large\bfseries}{\thesection}{0.8em}{}
\titlespacing*{\section}{0pt}{1.6ex plus 0.6ex minus 0.2ex}{0.9ex plus 0.2ex}
\titleformat{\subsection}{\raggedright\large\bfseries}{\thesubsection}{0.8em}{}
\titlespacing*{\subsection}{0pt}{1.3ex plus 0.5ex minus 0.2ex}{0.7ex plus 0.2ex}
\titleformat{\subsubsection}{\raggedright\normalsize\bfseries}{\thesubsubsection}{0.8em}{}
\titlespacing*{\subsubsection}{0pt}{1.0ex plus 0.4ex minus 0.1ex}{0.5ex plus 0.1ex}

\setlist[itemize]{
  leftmargin=2.4em,
  itemsep=0.25em,
  topsep=0.45em,
  parsep=0pt,
  partopsep=0pt
}
\setlist[enumerate]{
  leftmargin=2.6em,
  itemsep=0.25em,
  topsep=0.45em,
  parsep=0pt,
  partopsep=0pt
}

\tcbset{
  colback=yellow!8,
  colframe=red!55!black,
  boxrule=0.6pt,
  arc=1mm,
  left=1mm,
  right=1mm,
  top=0.8mm,
  bottom=0.8mm
}

\pagestyle{fancy}
\fancyhf{}
\fancyhead[L]{\nouppercase{\leftmark}}
\fancyhead[R]{机械制造基础}
\fancyfoot[C]{\thepage}
\renewcommand{\headrulewidth}{0.4pt}
\renewcommand{\footrulewidth}{0pt}
\renewcommand{\sectionmark}[1]{\markboth{\thesection\ #1}{}}
\setlength{\headheight}{15pt}

\setcounter{topnumber}{3}
\setcounter{bottomnumber}{2}
\setcounter{totalnumber}{5}
\renewcommand{\topfraction}{0.9}
\renewcommand{\bottomfraction}{0.8}
\renewcommand{\textfraction}{0.08}
\renewcommand{\floatpagefraction}{0.8}

\allowdisplaybreaks[2]
\clubpenalty=10000
\widowpenalty=10000
\displaywidowpenalty=10000
\emergencystretch=2em

\begin{document}
\maketitle
\newpage
\tableofcontents
\newpage

\begin{spacing}{1.22}
\RaggedRight
\section{机械制造技术基础课程笔记（按章节整理）}
\noindent\textbf{说明：}本笔记根据课件全文重排为“章-节”结构，保留原始知识点并修复断行。

\section{课程导学}
\subsection{总体要求与课程说明}
\begin{tcolorbox}[colback=yellow!6,colframe=red!55!black,title=重点,left=1mm,right=1mm]
\begin{enumerate}[leftmargin=2.2em,label=\arabic*.]
  \item 教学内容教学目的
  \item 考试形式考试内容
  \item 机械的历史与机械制造业的诞生
\end{enumerate}
\end{tcolorbox}

\section{第一章 绪 论}
\subsection{本章核心知识}
\begin{tcolorbox}[colback=yellow!6,colframe=red!55!black,title=重点,left=1mm,right=1mm]
\begin{enumerate}[leftmargin=2.2em,label=\arabic*.]
  \item \textbf{关键发明：现代机床的雏形}：1797年莫兹利发明的带移动刀架和导轨的机床，解决了精密加工难题，为机械制造提供了核心装备。
  \item \textbf{国民经济的支柱产业}：是国家创造力、竞争力和综合国力的重要体现，立国之本、兴国之器、强国之基，支撑着国民经济的
  \item \textbf{生产能力：规模领先}：多项核心产品产量位居世界第一，包括汽车、金属切削机床、发电设备等，展现了强大的制造规模优势。
  \item \textbf{自给率提升：突破瓶颈}：国内市场设备自给率实现跨越式增长，从不足60\%提升至超过85\%，关键领域国产化替代成效显著。
  \item \textbf{自主创新：技术突破}：在超超临界机组、核电、特高压等高端装备领域达到国际先进水平，成功打破国外长期技术垄断。
  \item \textbf{智能化与信息化}：向高度自动化、智能化、复合化方向发展，如智能制
\end{enumerate}
\end{tcolorbox}
\noindent\textbf{次重点与补充：}
\begin{itemize}[leftmargin=2.2em,label=--]
  \item 制造业萌芽：从手工到金属加工制造业起源于简单的木料加工，随着技术进步，逐渐过渡到对金属材料的精密加工，为机械制造奠定了物质基础。
  \item 第一次工业革命：动力的飞跃瓦特蒸汽机的发明提供了强大动力，推动了切削加工技术的巨大进步，标志着机械制造业的正式诞生。
  \item 发展。
  \item 各行各业的“装备部”为国民经济各部门提供现代化的技术装备，是其他
  \item 产业发展的物质基础。
  \item 科学技术转化的桥梁将科学技术转化为现实的生产力，推动产业升级和
  \item 经济发展。
  \item 历史表明：
  \item 没有制造能力的民族是没有竞争力的发展历程：从弱到强
  \item 从建国初期的薄弱基础起步，历经改革开放后的飞速发展，我国已成功跻身世界机械工业第一大国行列。
  \item 造、柔性制造系统。
  \item 高精度化超精密加工技术广泛应用，加工精度达到纳米级别，
  \item 满足高端制造需求。
  \item 新材料与新工艺增材制造（3D打印）等新技术不断涌现，推动材料
  \item 加工技术革新。
  \item 绿色与可持续发展绿色制造技术，注重环境保护和可持续发展，实
  \item 现节能减排目标。
\end{itemize}

\subsection{第一节 机械制造工业在国民经济中的地位与作用}
\begin{tcolorbox}[colback=yellow!6,colframe=red!55!black,title=重点,left=1mm,right=1mm]
\begin{enumerate}[leftmargin=2.2em,label=\arabic*.]
  \item 机械产品生产过程与机械制造过程
\end{enumerate}
\end{tcolorbox}
\noindent\textbf{次重点与补充：}
\begin{itemize}[leftmargin=2.2em,label=--]
  \item 机械制造工业在国民经济中的地位我国机械制造业的发展现状与趋势
\end{itemize}

\subsection{第二节 机械制造过程及机械制造系统}
\begin{tcolorbox}[colback=yellow!6,colframe=red!55!black,title=重点,left=1mm,right=1mm]
\begin{enumerate}[leftmargin=2.2em,label=\arabic*.]
  \item 生产过程三阶段涵盖从市场需求分析到最终产品交付的全过程，
  \item \textbf{制造方式分类}：- 减材制造：$\Delta G<0$，去除材料（如切削）
  \item \textbf{机械制造过程定义}：毛坯制造、机械加工、热处理、装配等直接改变零件形状、
  \item 核心要素：由机床、夹具、刀具和工件组成。
  \item 功能：机床提供动力，夹具装夹工件，刀具直接进行切削，四者协同完成加工。
  \item \textbf{机械制造系统}：系统视角：从系统论角度看，制造是一个输入输出系
  \item 转化过程：将生产要素（人、财、物、信息）转化为产品和效益。
\end{enumerate}
\end{tcolorbox}
\noindent\textbf{次重点与补充：}
\begin{itemize}[leftmargin=2.2em,label=--]
  \item 主要包括：决策阶段、设计研究阶段、制造阶段。
  \item - 等材制造：$\Delta G=0$，材料质量不变（如铸造）
  \item - 增材制造：$\Delta G>0$，逐层累加材料（如3D打印）
  \item 性能及位置关系的过程称为生产工艺过程。
  \item 机械加工工艺系统与机械制造系统机械加工工艺系统
  \item 统。
  \item 切削运动
\end{itemize}

\section{第二章 机械加工及设备的基础理论}
\subsection{本章核心知识}
\textit{（本节暂无可提取文本）}

\subsection{第一节 金属切削基本知识}
\begin{tcolorbox}[colback=yellow!6,colframe=red!55!black,title=重点,left=1mm,right=1mm]
\begin{enumerate}[leftmargin=2.2em,label=\arabic*.]
  \item \textbf{切削运动定义}：刀具和工件之间的相对运动，是实现切削的前提。
  \item \textbf{主运动}：速度最高、消耗功率最大的运动，如车削时工件的旋转。
  \item \textbf{进给运动}：将待切金属不断投入切削的运动，如车刀的移动。
  \item \textbf{合成切削运动}：主运动和进给运动的矢量和，决定了切削刃的运动轨迹。
  \item \textbf{切削速度 ($v_c$)}：切削刃相对于工件的主运动速度，单位为 $m/s$ 或 $m/min$；计算公式为 $v_c=\pi d_w n / 1000$。
  \item \textbf{进给量 ($f$)}：工件或刀具每转一周的相对位移，单位为 $mm/r$；进给速度可写为 $v_f=nf=n_zf_z$。
  \item \textbf{背吃刀量 ($a_p$)}：刀具切削刃与工件的接触长度在同时垂直于主运动和进给运动方向上的投影值，单位为 $mm$；外圆车削时 $a_p=(d_w-d_m)/2$。
  \item \textbf{车刀基本组成}：由刀柄（夹持部分）和刀头（切削部分）组成，其中刀头是承担切削任务的核心区域。
  \item \textbf{切削部分结构要素}：前刀面 ($A_\gamma$) 为切屑流过的刀面；后刀面 ($A_\alpha$) 为与工件过渡表面相对的刀面；副后刀面 ($A_\alpha'$) 为与工件已加工表面相对的刀面；主切削刃 ($S$) 为前刀面与后刀面的交线；副切削刃 ($S'$) 为前刀面与副后刀面的交线；刀尖是主、副切削刃的连接部分，可为圆弧或线段。
  \item \textbf{刀具角度的参考系}：基面 $P_r$ 为通过主切削刃上某一指定点并与该点切削速度方向相垂直的平面；切削平面 $P_s$ 为通过主切削刃上某一指定点、与主切削刃相切并垂直于该点基面的平面；正交平面 $P_o$ 为通过主切削刃上某一指定点，同时垂直于该点基面和切削平面的平面。
  \item \textbf{正交平面参考系}：由基面、切削平面和正交平面三个相互垂直的平面组成，是最常用的参考系。
  \item \textbf{刀具角度详解}：正交平面内角度包括前角 $\gamma_o$（前刀面与基面之间的夹角）、后角 $\alpha_o$（后刀面与切削平面之间的夹角）和楔角 $\beta_o$（前刀面与后刀面之间的夹角）。
  \item \textbf{切削平面内角度}：刃倾角 $\lambda_s$ 为主切削刃与基面之间的夹角，决定切屑的流出方向和刀尖强度。
  \item \textbf{基面内角度}：主偏角 $K_r$ 为主切削刃在基面上的投影与进给方向的夹角；副偏角 $K_r'$ 为副切削刃在基面上的投影与进给反方向的夹角；刀尖角 $\varepsilon_r$ 为主切削刃与副切削刃之间的夹角。
  \item \textbf{工作角度定义}：考虑合成切削运动和实际安装情况后的实际切削角度，区别于刀具的标注角度。
  \item \textbf{进给运动的影响}：进给运动会使刀具的工作前角增大、工作后角减小；在车削接近工件中心时，工作后角可能变为负值，导致工件被挤断；常用关系式为 $\gamma_{oe}=\gamma_o+\eta$，$\alpha_{oe}=\alpha_o-\eta$。
  \item \textbf{刀具安装对工作角度的影响}：当刀尖安装高于或低于工件中心时，会改变基面和切削平面的位置，从而直接影响刀具的工作角度；刀尖装高时工作前角增大、工作后角减小，刀尖装低时工作前角减小、工作后角增大。
  \item \textbf{切削层定义}：切削刃在一次走刀中从工件上切下的一层材料称为切削层，切削层的截面尺寸参数称为切削层参数。
  \item \textbf{公称厚度 ($h_D$)}：垂直于过渡表面的尺寸，$h_D=f\sin K_r$。
  \item \textbf{公称宽度 ($b_D$)}：沿过渡表面度量的切削层尺寸，$b_D=a_p/\sin K_r$。
  \item \textbf{公称横截面积 ($A_D$)}：切削层在基面上的投影面积，$A_D=h_D b_D=f a_p$。
\end{enumerate}
\end{tcolorbox}
\noindent\textbf{次重点与补充：}
\begin{itemize}[leftmargin=2.2em,label=--]
  \item 车削中主运动通常表现为工件旋转，合成切削运动轨迹直接决定已加工表面的形成。
  \item 切削用量三要素为切削速度 $v_c$、进给量 $f$ 和背吃刀量 $a_p$，其单位分别常取 $m/min$、$mm/r$ 和 $mm$。
  \item 对多刃刀具，进给速度常写为 $v_f=n_zf_z$；对车削，也常写为 $v_f=nf$。
  \item 正交平面 $P_o$ 是分析刀具标注角度时最常用的参考面。
  \item 正交平面内重点分析 $\gamma_o$、$\alpha_o$、$\beta_o$；切削平面内重点分析 $\lambda_s$；基面内重点分析 $K_r$、$K_r'$ 和 $\varepsilon_r$。
  \item 刃倾角 $\lambda_s$ 不仅影响切屑流向，也会影响刀尖强度和刀具受力状态。
  \item 切削层参数是后续分析切削变形、切削力和切削热的基础。
  \item 钻孔、端面铣、圆周铣等加工中同样需要分析切削用量三要素，但计算关系需结合具体刀具运动形式理解。
\end{itemize}

\subsection{第二节 金属切削机床的基本知识}
\begin{tcolorbox}[colback=yellow!6,colframe=red!55!black,title=重点,left=1mm,right=1mm]
\begin{enumerate}[leftmargin=2.2em,label=\arabic*.]
  \item \textbf{机床分类}：使用范围：通用机床、专门化机床、专用机床自动化程度：手动机床、机动机床、半自动机床、自动机床工作精度： 普通机床、精密机床、高精度机床质量和尺寸：仪表机床、中型机床、大型机床、重型机床
  \item \textbf{型号编制方法}：类代号：如“C”代表车床，用于区分机床类别。
  \item 组、系代号：用两位数字表示，进一步细分机床类型。
  \item \textbf{发生线的定义}：母线和导线统称为发生线，是由刀具的切削刃与工件的相
  \item \textbf{辅助运动}：除表面成形运动外的所有运动，用于实现快速移动、分度、
  \item 外联系传动链- 功能：联系动力源和执行件的传动链
  \item - 特点：不要求动力源和执行件之间有严格的传动比关系。
  \item \textbf{内联系传动链}：- 功能：联系两个执行件，形成复合成形运动的传动链- 特点：要求两个末端件之间有严格的传动比关系，直
\end{enumerate}
\end{tcolorbox}
\noindent\textbf{次重点与补充：}
\begin{itemize}[leftmargin=2.2em,label=--]
  \item 特性代号：如“K”代表数控，表示机床的通用特性。
  \item 主参数：表示机床规格的折算值，反映加工能力。
  \item 重大改进顺序号：如“A”、“B”，表示改进的次数。
  \item 工件表面的形成表面形成原理
  \item 任何工件表面都可看作是一条线（母线）沿着另一条线（导线）运动的轨迹
  \item 对运动获得的发生线的形成过程
  \item 通过刀具切削刃与工件的相对运动实现，机床的核心功能就是提供这些精确的相对运动。
  \item 机床的运动表面成形运动
  \item - 简单成形运动：单一旋转或直线运动，如车削外圆。
  \item - 复合成形运动：由多个简单运动组合，如车削螺纹的螺旋运动。
  \item 夹紧、松开等辅助动作。
  \item 机床的传动链传动链的定义
  \item 为了实现机床的运动，通过传动装置将动力源和执行件，或把执行件和执行件链接起来构成传动联系，构成这一联系的一系列传动件称为传动链。
  \item - 示例：外圆切削的传动链。
  \item 接影响加工精度。
  \item - 示例：车螺纹的传动链。
  \item 切屑的形成过程及变形区的划分
\end{itemize}

\subsection{第三节 金属切削过程}
\begin{tcolorbox}[colback=yellow!6,colframe=red!55!black,title=重点,left=1mm,right=1mm]
\begin{enumerate}[leftmargin=2.2em,label=\arabic*.]
  \item \textbf{第一变形区（剪切滑移区）}：切削层金属在刀具挤压下发生塑性变形，沿剪切面滑移形成切屑，是切削变形的主要区域。
  \item \textbf{第二变形区（纤维化区）}：切屑沿前刀面排出时，底层金属受强烈摩擦和挤压作用，发生进一步塑性变形并产生纤维化。
  \item \textbf{第三变形区（加工表面变形区）}：工件已加工表面受刃口挤压和后刀面摩擦作用，会产生弹塑性变形，并可能造成加工硬化和残余应力。
  \item \textbf{切屑的形成过程及变形区的划分}：OA、OB 和 OM 等均可视作等剪应力曲线，其中 OA 称为始滑移线，OM 称为终滑移线。
  \item \textbf{变形系数 $\Lambda_h$}：切屑厚度 $h_{ch}$ 与切削层厚度 $h_D$ 之比称为厚度变形系数 $\Lambda_{ha}$；切削层长度 $l_c$ 与切屑长度 $l_{ch}$ 之比称为长度变形系数 $\Lambda_{hl}$。$\Lambda_h$ 的值大于 1，且 $\Lambda_h$ 越大，切削变形越大；$\phi$ 增大时，$\Lambda_h$ 减小，切削变形减小。
  \item \textbf{相对滑移（剪应变）$\varepsilon$}：相对滑移图中，平行四边形 OHNM 发生剪切变形后变为平行四边形 OGPM，可用 $\varepsilon=\Delta s/\Delta y=\cot\phi+\tan(\phi-\gamma_o)$ 表示。
  \item \textbf{剪切角 $\phi$}：滑移剪切角 $\phi$ 与切削变形有密切关系，也常用来衡量切削变形程度。因为主应力方向与最大剪应力方向的夹角应为 $45^\circ$，即 $F_s$ 与合力 $F$ 的夹角应为 $45^\circ$。
  \item \textbf{剪切角变化规律}：前角 $\gamma_o$ 增大时，剪切角 $\phi$ 增大，切削变形减小；摩擦角 $\beta$ 增大时，剪切角 $\phi$ 减小，切削变形增大。
  \item \textbf{切屑类型}：带状切屑多见于塑性较好的材料，切削过程平稳、表面质量较好但不易断屑；节状切屑多见于中等塑性材料，切屑上有明显节状纹路，切削过程伴有轻微振动；单元切屑多见于塑性较低材料，切屑呈粒状，振动较大；崩碎切屑多见于脆性材料（如铸铁），切屑呈不规则碎块状，加工表面较粗糙。
  \item \textbf{积屑瘤}：切削钢和铝合金等塑性金属时，在切削速度不高而又能形成带状切屑的情况下，常在前刀面切削处粘着一块断面呈三角状的硬块，其硬度可达工件材料硬度的 2～3 倍，能够代替刀面和切削刃进行切削，这个硬块称为积屑瘤。
\end{enumerate}
\end{tcolorbox}
\noindent\textbf{次重点与补充：}
\begin{itemize}[leftmargin=2.2em,label=--]
  \item 一般切削速度下，第一变形区的宽度仅为 $0.02\sim0.2\,\mathrm{mm}$，因此常可近似用一个剪切面来表示。
  \item 剪切面与切削速度方向的夹角称为剪切角，记作 $\phi$。
  \item 变形系数越大，说明切屑压缩越严重，切削变形越明显。
  \item 带状切屑常见于低碳钢等塑性较好的材料，切削平稳、表面质量好，但不易断屑。
  \item 节状切屑常见于中等塑性材料，切削过程有轻微振动。
  \item 单元切屑多见于塑性较低材料，切削过程振动较大，会影响加工精度。
  \item 崩碎切屑多见于脆性材料，容易导致加工表面粗糙。
  \item 积屑瘤形成的典型条件是：切削塑性材料、前刀面存在强烈摩擦、局部压力较高且温度适当。
  \item 积屑瘤会使刀具实际前角变化，导致切削厚度波动并使加工表面粗糙度增大；它对刀具寿命和切削过程既可能有暂时的保护作用，也会带来不稳定影响。
  \item 精加工时应尽量防止积屑瘤产生，可采取避开易形成积屑瘤的切削速度范围、合理使用切削液、增大刀具前角或降低材料塑性的措施。
  \item 切削力、切削热和切削温度的分析，都建立在对切削变形过程的理解基础上。
\end{itemize}

\subsection{第四节 切削力、切削热与切削温度}
\begin{tcolorbox}[colback=yellow!6,colframe=red!55!black,title=重点,left=1mm,right=1mm]
\begin{enumerate}[leftmargin=2.2em,label=\arabic*.]
  \item \textbf{切削力的来源}：- 克服被加工材料弹性变形的抗力- 克服被加工材料塑性变形的抗力- 克服切屑与刀具前刀面的摩擦力- 克服刀具后刀面与工件表面的摩擦力切削力的分解 (三个相互垂直分力)
  \item 切削用量的影响- 背吃刀量ap和进给量f ：
  \item \textbf{- 切削速度vc ：}：0 40 80 120 160 200 240 280
  \item \textbf{vc(m/min)}：切削塑性金属时，在中速和高速下，切削力一般随着vc的增大而减小。这主要是因为vc增大，使切削温度提高，摩擦因数减小，从而使变形系数减小在低速范围内，由于存在着积屑瘤，切削速度对切削力的影响有着特殊的规律
  \item \textbf{刀具几何参数的影响}：前角 $\gamma_o$ 对切削塑性材料时的切削力影响较大；切削脆性材料时，由于切削变形较小，$\gamma_o$ 对切削力的影响不显著。
  \item \textbf{主偏角 $k_r$}：$k_r$ 对主切削力 $F_c$ 的影响不大，但对背向力 $F_p$ 和进给力 $F_f$ 的影响较大；工艺系统刚性较差时，可通过改变主偏角来减小背向力。
  \item \textbf{刃倾角 $\lambda_s$}：改变刃倾角会影响切屑在前刀面上的流动方向，从而使切削合力方向发生变化；增大 $\lambda_s$ 时，$F_p$ 减小、$F_f$ 增大，在 $\lambda_s$ 于 $-45^\circ\sim10^\circ$ 范围变化时，$F_c$ 基本不变。
  \item \textbf{刀尖圆弧半径 $r_\varepsilon$}：$r_\varepsilon$ 对 $F_p$、$F_f$ 的影响较大，对 $F_c$ 的影响较小。
  \item \textbf{刀具磨损的影响}：前刀面磨损使刀具实际前角增大，切削力减小。
  \item \textbf{切削液的影响}：切削液的作用：润滑、冷却、清洗、除油
  \item \textbf{以冷却作用为主的水溶液对切削力的影响较小}：使用润滑性能好的切削液，能有效减少摩擦，显著降低切削力
  \item \textbf{- 切屑 ：带走绝大部分热量。}：- 刀具 ：导致刀具磨损和破损的主因。
  \item \textbf{- 工件 ：引起工件热变形，影响精度。}：- 周围介质：通过辐射、对流散发，占比极少。
  \item \textbf{切削温度的定义}：- 切削温度一般是指切屑与前刀面接触区域的平均温度
  \item \textbf{切削温度的测量}：- 自然热电偶 ：切屑与前刀面接触区的平均温度。
  \item \textbf{淬火}：工件材料：合金结构钢的强度普遍高于45钢，而导热系数又一般均低于45钢。所以切削合金结构钢时的切削温度一般均高于切削45钢时的切削温度
  \item \textbf{刀具参数对切削温度的影响}：前角 $\gamma_o$ 增大时，切削变形减小，切削温度一般下降；主偏角 $k_r$ 也会影响切削温度，减小 $k_r$ 时，切削刃工作长度和刀尖角增大，散热条件改善，切削温度通常下降。
  \item 刀具磨损：刀具磨损后，切削刃变钝，刃区前方的挤压作用增大，使切削区的金属的塑性变形增加。同时，磨损后的刀具后角变成零度，使工件与刀具的摩擦加大，两者均使切削热的产生增加。所以，刀具磨损是影响切削温度的重要因素。
  \item \textbf{vc＝71m/min}：0 0.2 0.4 0.6 0.8 VB(mm)
  \item 切削液：切削液对降低切削温度、减少刀具磨损和提高已加工表面质量有明显的效果，在切削加工中应用很广。
\end{enumerate}
\end{tcolorbox}
\noindent\textbf{次重点与补充：}
\begin{itemize}[leftmargin=2.2em,label=--]
  \item 切削力可分解为主切削力 $F_c$、进给力 $F_f$ 和背向力 $F_p$，其中 $F_c$ 沿主运动方向、数值最大、消耗功率最多。
  \item 切削力来源可概括为材料弹塑性变形抗力以及切屑、工件与刀具前后刀面的摩擦阻力。
  \item 切削力经验公式表明：切削用量对切削力的影响程度通常满足 $a_p > f > v_c$。
  \item - $F_c = C_{F_c} a_p^{x_{F_c}} f^{y_{F_c}} v_c^{n_{F_c}} K_{F_c}$
  \item - $F_p = C_{F_p} a_p^{x_{F_p}} f^{y_{F_p}} v_c^{n_{F_p}} K_{F_p}$
  \item - $F_f = C_{F_f} a_p^{x_{F_f}} f^{y_{F_f}} v_c^{n_{F_f}} K_{F_f}$
  \item 切削力的测量
  \item 影响切削力的因素工件材料的影响
  \item - 工件材料的强度越高，切削力越大。
  \item - 切削脆性材料时，被切材料的塑性变形及它与前刀面的摩擦都比较小，故其切削力相对较小。
  \item $a_p$ 和 $f$ 增大都会使切削力增大，但两者的影响程度不同。
  \item $a_p$ 增大时，变形系数 $\Lambda_h$ 基本不变，切削力近似成正比增大。
  \item $f$ 增大时，$\Lambda_h$ 有所下降，因此切削力不成正比增大。
  \item 切削铸铁等脆性材料时，被切材料的塑性变形及与前刀面的摩擦均较小，因此 $v_c$ 对切削力往往没有显著影响。
  \item 当前角 $\gamma_o$ 增大时，$\Lambda_h$ 减小，切削力下降。
  \item F(kN)
  \item FcFf
  \item \textbf{负倒棱}：当 $b\gamma_1<l_f$ 时，切削力虽有所增加，但增大幅度不大；当 $b\gamma_1>l_f$ 时，相当于用负前角为 $\gamma_{o1}$ 的车刀进行切削，切削力显著增加。
  \item \textbf{主偏角的进一步影响}：主偏角 $k_r$ 增大时，背向力 $F_p$ 减小、进给力 $F_f$ 增大，而 $F_c$ 基本不变。
  \item \textbf{刀尖圆弧半径的进一步影响}：当 $r_\varepsilon$ 增大时，平均主偏角减小，因此 $F_p$ 减小、$F_f$ 增大。
  \item 后刀面磨损，刀具与工件的摩擦增大，切削力增大。
  \item 前后刀面同时磨损时，切削力先减小，后逐渐增大。
  \item 切削热的产生与传导切削热的产生来源
  \item - 能量转化：切削层金属发生弹性和塑性变形消耗能量转化为热- 摩擦生热：刀具与切屑、工件之间的摩擦功转化为热- Q=Fcvc
  \item 切削用量三要素对切削热的影响 ap > vc > f切削热的传导途径
  \item 传 导 途 径 干 车 削 钻 削切屑 50\%～86\% 28\%工件 40\%～10\% 52\%刀具 9\%～3\% 15\%
  \item 周围介质 1\% 5\%切削温度的测量方法
  \item - 人工热电偶 ：距前刀面一定距离处某点的温度，不能直接测出前刀面上的温度。
  \item 毫伏表工件 车床主轴尾部
  \item 铜销铜顶尖(与支架绝缘)
  \item 导线卡盘
  \item 绝缘层顶尖
  \item (a)测刀具 (b)测工件刀具 工件
  \item 影响切削温度的主要因素关键影响因素分析
  \item 工件材料：工件材料的强度、硬度越高，切削力越大，切削时消耗的功也越多，产生的切削热也越多，切削温度也就越高
  \item 10 20 30 40 50 60 70 80 90 100 110 120 130theta(℃)
  \item 调质正火
  \item 45 正火W18Cr4
  \item V5CrNiM
  \item o38CrMoAl
  \item A工件材料：
  \item 不锈钢1Cr18Ni9Ti和高温合金GH13不但导热系数低，而且在高温下仍能保持较高的强度和硬度。所以切削这种类型的材料时，切削温度比切削其他材料要高得多脆性金属的抗拉强度和延伸率都较小，切削过程中切削区的塑性变形很小，切屑呈崩碎状或脆性带状，与前刀面的摩擦也很小，所以产生的切削热较少，切削温度一般比切削钢料时低
  \item GH131 45 钢(正火)
  \item HT20～401Cr18Ni9Ti
  \item 切削用量：vc ＞f ＞ap随着vc的增大，变形热与摩擦热增多，在很短的时间内，切屑底层的切削热来不及向切屑和刀具内部传导，而积聚在切屑底层，从而使切削温度显著升高。
  \item 随着 f 的增加，一方面金属切除率增多，切削温度升高;另一方面单位切削力和单位切削功率减小，切除单位体积金属所产生的热量减少。此外，当f增大后，切屑变厚，由切屑带走的热量增多，故切削温度上升不显著。
  \item ap 增大后，产生的热量虽成比例增多，但因切削刃参加工作的长度也成正比增长，改善了散热条件，所以切削温度升高得不明显。
  \item 切削温度经验关系可概括为 $\theta=C_\theta v_c^x f^y a_p^z$，其中 $\theta$ 表示刀具与切屑接触区的平均温度。
  \item 一般情况下，切削用量对切削温度的影响大小规律为 $v_c > f > a_p$。
  \item 当前角 $\gamma_o$ 增大时，切削变形减小，切削力减小，切削温度下降。
  \item 减小主偏角 $k_r$ 时，切削刃工作长度和刀尖角增大，散热条件改善，因此切削温度也会下降。
  \item vc＝2.25m/svc＝1.35m/s
  \item vc＝117m/minvc＝94m/min
  \item 切削液对切削温度的影响，与切削液的导热性能、比热、流量、浇注方式以及本身的温度有很大关系。
  \item 从导热性能来看，水基切削液＞乳化液＞油类切削液刀具磨损的形式
\end{itemize}

\subsection{第五节 刀具磨损与刀具寿命}
\begin{tcolorbox}[colback=yellow!6,colframe=red!55!black,title=重点,left=1mm,right=1mm]
\begin{enumerate}[leftmargin=2.2em,label=\arabic*.]
  \item \textbf{前刀面磨损（月牙洼磨损）}：切削塑性材料时，在刀具前刀面靠近切削刃处形成月牙形磨损区域，由剧烈摩擦和高温导致。
  \item \textbf{后刀面磨损}：刀具后刀面与工件过渡表面摩擦形成磨损带，是最常见
  \item 形式，直接影响工件精度和粗糙度。
  \item \textbf{前、后刀面同时磨损}：切削塑性材料时，切削层公称厚度hd适中 。
  \item \textbf{磨料磨损}：也称机械磨损，由于切屑或工件的摩擦面上有一些微小的硬质点，能在刀具表面刻划出沟纹，这就是磨料磨损。
  \item \textbf{扩散磨损}：刀具材料和工件材料在高温下化学元素相互扩散而造成的磨损。在高温下(900\textasciitilde{}1000°C)刀具材料中的Ti、W、Co等元素会扩散到切屑或工件材料中去，而工件材料中的Fe元素也会扩散到刀具表层里，这样改变了硬质合金刀具的化学成分，使表层硬度变低，从而加剧刀具磨损
  \item \textbf{氧化磨损}：当切削温度达700\textasciitilde{}800°C时，空气中的氧与硬质合金中的钴及碳化钨、碳化钛等发生氧化作用，产生较软的氧化物(如CoO、WO2、TiO2)被切屑或工件摩擦掉而形成
  \item \textbf{磨钝标准}：刀具磨损到一定限度就不能再继续使用，这个磨损限度
  \item \textbf{称为磨钝标准。}：国际标准ISO统一规定，刀具磨钝标准是指后刀面磨损带中间部分平均磨损量允许达到的最大值，以VB表示。
  \item \textbf{磨钝标准的制定}：既能发挥刀具的切削能力，又要考虑保证工件的加工质量。
  \item \textbf{粗加工时，磨钝标准取较大值；}：工艺系统刚性较差时，磨钝标准取较小值；
  \item \textbf{0.4～0.5}：碳素钢粗车 0.6～0.8铸铁件粗车 0.8～1.2
  \item \textbf{刀具寿命的定义}：刀具寿命是指新刃磨的刀具从开始切削一直到磨损量达到磨钝标准时的切削时间，用符号T表示，单位为s(或min)
\end{enumerate}
\end{tcolorbox}
\noindent\textbf{次重点与补充：}
\begin{itemize}[leftmargin=2.2em,label=--]
  \item 月牙洼磨损一般出现在前刀面靠近主切削刃的位置；后刀面磨损常以 $V_B$ 表示。
  \item 刀具磨损形式常结合主切削刃、副切削刃及剖面示意图理解。
  \item 刀具磨损的原因主要包括磨料磨损、粘结磨损、扩散磨损和氧化磨损。
  \item 磨料磨损在各种切削速度下都存在，但对低速切削刀具(如拉刀、板牙等)磨料磨损是主要原因粘结磨损
  \item 粘结磨损也称冷焊磨损。切屑或工件的表面与刀具表面之间发生粘结现象。由于有相对运动，刀具上的微粒被对方带走而造成磨损。
  \item 粘结磨损与切削温度有关，也与刀具及工件两者的化学成分有关(元素的亲和作用)。
  \item 的磨损称为氧化磨损。
  \item 刀具磨损过程及磨钝标准刀具的磨损过程
  \item 1. 初期磨损阶段 该阶段磨损曲线的斜率较大，这意味着刀具磨损很快
  \item 2. 正常磨损阶段 经过初期磨损后，刀具的后刀面上被磨出一条狭窄的棱面，压强减小。同时刀具的表面已经被磨平，磨损量的增加减缓并稳定下来，刀具进入正常磨损阶段
  \item 3. 急剧磨损阶段 刀具经过正常磨损阶段后，切削刃明显变钝，引起切削力、切削温度迅速增大。这时进入急剧磨损阶段，这一阶段磨损曲线斜率很大，表现为刀具磨损速度很快
  \item 0 40 80 120 160 200切削时间 t(min)
  \item 0.40.3
  \item 0.20.1
  \item 初期磨损 正常磨损 急剧磨损VB(mm)
  \item 因为刀具的后刀面上都有磨损，它对加工精度和切削力的影响比前刀面磨损显著，且易于测量，因此在刀具管
  \item 理中多按后刀面磨损尺寸来制订磨钝标准。
  \item 后刀面磨损增大时，实际后角会减小，因此精加工时磨钝标准应取较小值；
  \item 切削难加工材料时，磨钝标准也取较小值。
  \item 硬质合金车刀的磨钝标准加 工 条 件 后刀面的磨钝标准
  \item VB/mm精车 0.1～0.3
  \item 合金钢粗车，粗车刚性较差的工件
  \item 钢及铸铁件大件低速粗车
  \item 1.0～1.5刀具寿命
  \item 切削用量与刀具寿命的关系cp
  \item Tmmm
  \item CTv f a
  \item 5 2.25 0.75TCT v f a
  \item 在切削用量三要素中，切削速vc对刀具耐用度的影响最大，其次为进给量f ，背吃刀量ap的影响最小，这与三者对切削温度的影响顺序是一致的
  \item vc ＞f ＞ap硬质合金外圆车刀切削碳钢时
  \item 刀具合理寿命的选择最大生产效率使用寿命：
  \item 经济寿命：
  \item 刀具结构复杂、制造和磨刀费用高时，刀具寿命规定高些
  \item 多刀车床上的车刀，组合机床上的刀具，自动线上的刀具，刀具寿命应规定高些
  \item 某工序成为生产瓶颈时，刀具寿命要规定的低些 某工序单位时间的生产成本较高时，刀具寿命规定的低些
  \item 加工大型工件时，刀具寿命应规定的高些衡量切削加工性的指标
\end{itemize}

\section{第三章 切削条件的合理选择及刀具的选择}
\subsection{本章核心知识}
\textit{（本节暂无可提取文本）}

\subsection{第一节 工件材料的切削加工性}
\begin{tcolorbox}[colback=yellow!6,colframe=red!55!black,title=重点,left=1mm,right=1mm]
\begin{enumerate}[leftmargin=2.2em,label=\arabic*.]
  \item \textbf{定义}：材料的切削加工性是指对某种材料进行切削加工的
  \item 难易程度衡量指标
  \item - 在相同的切削条件下，刀具有较长的寿命;或在一定的刀具寿命下，能够采用较高的切削速度。
  \item \textbf{v60j}：Kr ——某种材料相对加工性能指标；
  \item \textbf{体加工性好}：总结：材料的切削加工性是多种因素综合作用的结果，需综合考虑物理性能、成分及组织结构。
  \item \textbf{刀尖应力大，易崩刀}：切屑与前刀面接触长度短，导致刀尖附近应力集中，容易产生塑性变形或崩刃。
  \item \textbf{导热性差，切削温度高}：与高锰钢类似，热量积聚在切削区，加速刀具磨损
  \item \textbf{切削对策}：通常选用韧性好的YG类硬质合金，采用较大的前角、较小的主偏角，以及较低的切削速度和较大的进给量。
  \item \textbf{采用适当的热处理}：低碳钢正火提高硬度；高碳钢球化退火降低硬度；铸铁退火消除白口组织。
  \item \textbf{采用新的切削加工技术}：如加热、低温、振动切削等。加热切削可降低材料抗剪强度，减小切削力。
\end{enumerate}
\end{tcolorbox}
\noindent\textbf{次重点与补充：}
\begin{itemize}[leftmargin=2.2em,label=--]
  \item - 在相同的切削条件下，切削力或切削功率小，切削温度低。
  \item - 容易获得良好的表面加工质量。
  \item - 容易控制切屑的形状或容易断屑加工性等级 工件材料分类 相对
  \item 加工性指标Kr 代表性材料1 很容易切削的材料 一般有色金属 >3.0 铜铅合金、铝镁合金、铝青铜容易切削的材料
  \item 易切钢 >2.5\textasciitilde{}3.0 退火15Cr、自动机钢3 较易切钢 >1.6\textasciitilde{}2.5 正火30钢普通材料
  \item 一般钢、铸铁 >1.0\textasciitilde{}1.6 45 钢、灰铸铁、结构钢5 稍难切削的材料 >0.65\textasciitilde{}1.0 调质2Cr13、85钢难切削的材料
  \item 较难切削的材料 >0.5\textasciitilde{}0.65 调质45Cr、调质65Mn7 难切削的材料 >0.15\textasciitilde{}0.5 1Cr18Ni9Ti、调质50CrV、某些钛合金
  \item 8 很难切削的材料 <0.15 铸造镍基高温合金、某些钛合金Kr= v60
  \item v60 ——某种材料其刀具寿命为60min时的切削速度；
  \item v60j ——切削45钢，刀具寿命为60min时的切削速度。
  \item 影响材料切削加工性的主要因素物理力学性能
  \item - 强度硬度：越高加工性越差。
  \item - 塑性韧性：越强变形大，断屑难。
  \item - 导热性：越好越利于散热，加工性好。
  \item 化学成分- 碳钢：中碳钢加工性最好，低
  \item 碳钢塑性大。
  \item - 合金元素：锰、硅、镍等提高硬度，降低加工性。
  \item - 铸铁：碳、硅等促进石墨化，改善加工性。
  \item 金相组织- 钢：珠光体加工性好，渗碳体
  \item 、马氏体差- 铸铁：铁素体基体比珠光体基
  \item 难加工材料的切削加工性 - 高锰钢加工硬化严重
  \item 切削过程中，奥氏体组织转变为细晶粒马氏体，导致表
  \item 面硬度急剧增加，造成切削困难。
  \item 热导率低热导率仅为45钢的1/4，切削
  \item 热量不易散发，导致切削温度很高，加剧刀具磨损。
  \item 韧性大韧度约为45钢的8倍，伸长率
  \item 大，切削变形严重，切削力大，且切屑不易折断。
  \item 总结：高锰钢的三大特性使其切削力大、温度高、刀具磨损快，是切削加工中的一大挑战。
  \item 难加工材料的切削加工性 – 不锈钢加工硬化严重，易粘刀
  \item 不锈钢中的铬、镍元素使其加工硬化倾向明显，且切屑容易粘附在前刀面上。
  \item 。
  \item 切屑为锯齿形，易振动切削过程不平稳，容易产生振动，影响加工质量和
  \item 刀具寿命。
  \item 改善材料切削加工性的基本方法添加化学元素
  \item 钢材中添加硫、铅等破坏铁素体连续性；
  \item 铸铁中加入铝、铜等促进石墨化元素。
  \item 图3-1 切削不锈钢时切削温度与切削力关系结论：温度升高，切削力显著下降刀具材料应具备的性能
\end{itemize}

\subsection{第二节 刀具材料}
\begin{tcolorbox}[colback=yellow!6,colframe=red!55!black,title=重点,left=1mm,right=1mm]
\begin{enumerate}[leftmargin=2.2em,label=\arabic*.]
  \item \textbf{高的硬度}：刀具材料的硬度必须高于工件材料，常温下硬度应在60HRC以上
  \item 高的耐磨性能抵抗切削过程中的磨损，保持
  \item \textbf{强度和韧性}：以承受切削力和冲击，防止刀具崩刃或折断，确保加工安全。
  \item \textbf{高的耐热性}：在高温下仍能保持较高的硬度、耐磨性和强度，即高温硬度。
  \item \textbf{导热性与工艺性}：导热性好利于散热，工艺性好便于刀具的制造、刃磨和维护。
  \item \textbf{普通高速钢}：如W18Cr4V、W6Mo5Cr4V2，耐热性约600°C，能满足一般工程材料的常规加工需求。
  \item \textbf{高性能高速钢}：通过调整化学成分，切削性能显著提高，主要用于加工高温合金、钛合金、不锈钢等难加工材料。
  \item 丝锥、成形铣刀及齿轮刀具等核心价值：在不改变高速钢基体韧性的前提下，极大地提高表面硬度和耐磨性，实现刀具寿命的跨越式提升。
  \item \textbf{核心性能}：- 硬度高达89\textasciitilde{}94HRA，耐热性800\textasciitilde{}1000°C- 切削速度比高速钢高4\textasciitilde{}10倍，寿命长- 缺点：抗弯强度低、韧性差，不耐冲击
  \item \textbf{加工长切屑的钢件}：P类 (YT-WC、Co、TiC类)
  \item \textbf{通用合金，钢或铸铁}：M类 (YW-WC、Co、TiC、TaC或NbC类)
  \item \textbf{加工短切屑的铸铁}：N类(WC、Co、TaC、 NbC或CrC类)
  \item \textbf{有色金属和非金属}：S类(WC、Co、TaC、 NbC或TiC类)
  \item \textbf{耐热合金和优质合金}：H类(WC、Co、TaC、 NbC或TiC类)
  \item \textbf{淬硬钢、冷硬铸铁等}：应用总结：硬质合金凭借其优异的红硬性和耐磨性，广泛应用于切削刀具、模具制造等领域。正确选择合金类别（如P类用于钢件，K类用于铸铁）是确保加工效率和质量的关键。
  \item \textbf{牌号特点}：- 数字越大：钴含量越高，韧性越好，适合粗加工- 数字越小：耐磨性和耐热性越好，适合精加工。
  \item \textbf{应用示例}：- P01/P10：钢件的精加工/半精加工- K01/K10：铸铁的精加工/半精加工- M10/M20：钢或铸铁的半精加工
  \item \textbf{金刚石刀具}：目前最硬的物质(10000HV)，用于切削有色金属、非金属材料，寿命极高。
  \item \textbf{立方氮化硼 (CBN)}：硬度仅次于金刚石，耐热性高(1300\textasciitilde{}1500°C)。
  \item \textbf{高精度与抗振性}：热变形小，抗振性能优异，保证加工精度。
  \item \textbf{高互换性}：便于快速换刀，适应自动化生产节奏。
  \item \textbf{稳定性能与寿命}：切削性能稳定，长寿命降低生产成本。
  \item \textbf{主要数控刀具材料}：金刚石刀具(PCD) 立方氮化硼(PCBN) 陶瓷刀具 涂层刀具 超细晶粒硬质合金 粉末冶金高速钢
  \item \textbf{物理性能匹配}：考虑熔点、导热系数等，导热性好的刀具适合
  \item \textbf{化学性能匹配}：避免刀具与工件材料发生强烈的化学反应和扩
  \item \textbf{涂层刀具特点}：- 提高表面硬度，减少磨损摩擦化学/热屏障，
  \item \textbf{减少扩散和化学反应}：- 显著提升切削速度、提高加工精度，提高刀
  \item \textbf{具寿命}：- 支持干切削，符合绿色加工- 提高刀具表面韧度，减少断裂，
  \item \textbf{纳米涂层技术}：通过涂覆上百层纳米级材料，显著提高刀具的高温硬度和耐磨性。
  \item \textbf{常用涂层材料}：TiN（氮化钛）呈金黄色，应用广泛，适用于加工低合金钢和不锈钢；TiCN（碳氮化钛）硬度更高、呈灰色，适用于各种工件材料。
  \item CBN (立方氮化硼)：金黄色，可进行淬火钢工件批量加工切削液的作用
\end{enumerate}
\end{tcolorbox}
\noindent\textbf{次重点与补充：}
\begin{itemize}[leftmargin=2.2em,label=--]
  \item 。
  \item 刀具的锋利，延长使用寿命。
  \item 高速钢 - 普通与高性能概述
  \item 高速钢是加入了W、Mo、Cr、V等合金元素的高合金工具钢，具有较好的抗弯强度和韧性，可制造刃形复杂的刀具
  \item - 钴高速钢 (M42)：综合性能好，硬度高，但价格昂贵。
  \item - 铝高速钢 (W6Mo5Cr4V2Al)：我国独创，性能优异且价格低廉。
  \item - 高钒高速钢：耐磨性好，用于切削高强度钢。
  \item 高速钢 - 涂层高速钢先进工艺
  \item 采用 PVD（物理气相沉积）技术，在刀具表面涂覆一层厚度约为 $2\sim5\,\mu\text{m}$ 的 TiN、TiC、TiCN 等硬膜。性能优势：
  \item 表面硬度可达2200HV，耐磨性大幅提高。刀具使用寿命可提高3\textasciitilde{}7倍，切削效率显著提升。
  \item 广泛应用已广泛应用于形状复杂的钻头、
  \item 硬质合金 - 概述与分类材料概述
  \item 由金属碳化物粉末（如WC、TiC）和金属粘结剂（如Co）经粉末冶金制成，是当今最主要的刀具材料之一。
  \item 分类标准 (GB/T 18376.1-2008)
  \item K类 (YG-WC、Co类)
  \item 硬质合金 - 牌号选择与应用选择原则
  \item 主要根据工件材料和切削加工类型（粗加工/精加工）
  \item 来确定。
  \item 其他材料陶瓷刀具
  \item 具有极高硬度和耐磨性，耐热性可达1200°C以上。
  \item 但抗弯强度低、韧性差，不耐冲击。
  \item 主要用于高速精加工。
  \item 但不适合切削钢铁，热稳定性较差。
  \item 主要用于加工淬硬钢、高温合金等难加工材料，可实现“以车代磨”。
  \item 总结：这三类材料各有所长，陶瓷适合高速精加工，金刚石是硬度之王，CBN则是硬切削的理想选择。
  \item 数控加工刀具材料核心特点
  \item 高效 - 高精度 - 高可靠性 - 专用化高强度与刚性
  \item 承受高速、大负荷切削，确保加工稳定性。
  \item 刀具材料与工件材料的合理匹配力学性能匹配
  \item 刀具硬度应高于工件硬度，高温力学性能优异的刀具适合高速切削。
  \item 加工导热性差的工件。
  \item 散，减少化学磨损。
  \item 常用数控刀具材料匹配示例工件材料 推荐刀具材料
  \item 铝合金 铝高速钢、K类硬质合金普通钢 钴高速钢、硬质合金、金属陶瓷高强度钢/高温合金 超细晶粒硬质合金、金属陶瓷钛合金 聚晶金刚石(PCD)、涂层硬质合金复合材料 聚晶金刚石(PCD)、涂层硬质合金涂层刀具材料
  \item TiAlN (氮铝化钛)：耐高温，灰/黑色，适用于硬钢、钛合金高速切削。
  \item AlCrN (氮化铬铝)：高温稳定性好，可干切削，适用于合金钢、铸铁等。
  \item DLC (类金刚石)：硬度高，摩擦因数低，适用于加工铝材、石墨等。
  \item 总结：涂层技术是提升刀具性能的有效手段，通过在刀具表面涂覆薄而硬的涂层，可极大提高硬度、耐磨性和耐热性。
  \item TiAlSiN (氮化钛铝硅)：硬度高，灰/黑色，可干切削。
\end{itemize}

\subsection{第三节 切削液}
\begin{tcolorbox}[colback=yellow!6,colframe=red!55!black,title=重点,left=1mm,right=1mm]
\begin{enumerate}[leftmargin=2.2em,label=\arabic*.]
  \item \textbf{冷却作用}：通过热传导、对流和汽化带走切削热，降低切削温度，从而提高刀具寿命和加工质量。水溶液的冷却效
  \item 润滑作用减少切屑、工件与刀具之间的摩擦。
  \item \textbf{清洗作用}：冲走切削区的切屑和磨屑，防止划伤已加工表面和机床导轨，保持加
  \item 总结1. 核心作用：
  \item 乳化液：油水乳化，性能均衡，应用广泛。
  \item \textbf{选用原则}：工件材料：铸铁一般不用，铜铝忌含硫。
  \item \textbf{切削液的种类和选用}：序号 名称 组成 主 要 用 途1 水溶液 硝酸钠、碳酸钠等溶于水的溶液，用100\textasciitilde{}200 倍的水
  \item \textbf{4 其他}：液态的(CO2 主要用于冷却二硫化钼+硬脂酸+石蜡做成蜡笔，涂于刀具表面 攻螺纹
\end{enumerate}
\end{tcolorbox}
\noindent\textbf{次重点与补充：}
\begin{itemize}[leftmargin=2.2em,label=--]
  \item 果最好。
  \item 通过添加油性添加剂和极压添加剂（含S、P、Cl）来改善边界润滑效果。
  \item 工环境清洁。
  \item 防锈作用在金属表面形成保护膜，防止工件
  \item 和机床在加工过程中或加工后生锈腐蚀。
  \item 冷却与润滑是最关键的功能，直接决定切削力、温度及刀具寿命。
  \item 2. 辅助保障：
  \item 清洗与防锈确保了加工环境的清洁和设备工件的长期稳定性。
  \item 切削液的种类及选用切削液的种类
  \item 水溶液：冷却清洗性能好，主要用于磨削。
  \item 切削油：矿物油为主，润滑好，用于低速精密加工。
  \item 其他：液态CO2、固体润滑剂等。
  \item 刀具材料：硬质合金刀具需连续充分供应。
  \item 加工要求：粗加工侧重冷却，精加工侧重润滑切削液种类与选用指南
  \item 稀释而成 磨削2 乳化液
  \item 1)少量矿物油，主要为表面活性剂的乳化油，用40\textasciitilde{}80倍的水稀释而成，冷却和清洗性能好 车削、钻孔2)以矿物油为主，少量表面活性剂的乳化油，用10\textasciitilde{}20 倍的水稀释而成，冷却和清洗性能好 车削、攻螺纹3)在乳化液中加入极压添加剂 高速车削、钻削3 切削油
  \item 1) 矿物油(L-AN15 戒L-AN32 全损耗系统用油)单独使用 滚齿、插齿
  \item 2)矿物油加植物油戒动物油形成混合油，润滑性能好 精密螺纹车削3)矿物油戒混合油中加入极压添加剂形成极压油 高速滚齿、插齿、车螺纹
  \item 新型切削液
\end{itemize}

\subsection{第三节 新型切削液}
\begin{tcolorbox}[colback=yellow!6,colframe=red!55!black,title=重点,left=1mm,right=1mm]
\begin{enumerate}[leftmargin=2.2em,label=\arabic*.]
  \item \textbf{液滴直径小于 $0.1\,\mu\text{m}$}：- 兼具润滑性与冷却清洗性 - 无异味、无毒、稳定性好
  \item \textbf{- 注重环保可降解特性}：- 减少废液排放，绿色制造总结：新型切削液正朝着高性能、环保、长寿命的方向发展，是实现绿色制造的关键一环。
\end{enumerate}
\end{tcolorbox}
\noindent\textbf{次重点与补充：}
\begin{itemize}[leftmargin=2.2em,label=--]
  \item 微乳切削液- 介于乳化油和合成液之间，
  \item 新型水基切削液- 添加油性和极压添加剂
  \item - 显著改善润滑和防锈性能- 水基切削液的重大突破
  \item 发展趋势- 研制高性能、长寿命切削液
  \item 前角的功用及选择
\end{itemize}

\subsection{第四节 刀具合理几何参数的选择}
\begin{tcolorbox}[colback=yellow!6,colframe=red!55!black,title=重点,left=1mm,right=1mm]
\begin{enumerate}[leftmargin=2.2em,label=\arabic*.]
  \item \textbf{前角的功用}：增大前角：减小切屑变形和摩擦，降低切削力与温度，抑制积屑瘤，改善表面质量。
  \item \textbf{前角的选择原则}：- 工件材料：软、塑性大取大值；硬、脆性取小值。
  \item \textbf{后角的选择原则}：1. 切削层厚度：粗加工（厚）取小值；精加工（薄）
  \item \textbf{取大值。}：2. 工艺系统刚性：刚性差或尺寸精度要求高时，取较
  \item \textbf{小后角。}：3. 工件材料：材料强度、硬度越大，后角应取越小值
  \item 刃倾角的主要功用- 影响切屑流向：正值使切屑流向待加工表面
  \item - 影响刀刃强度：负值可保护刀尖，增加强度- 影响切削力比值：负值会增大背向力
  \item \textbf{刃倾角选择原则}：- 精加工：取正值 (0°\textasciitilde{}+5°)，防止划伤已加工表面- 粗加工：取负值 (0°\textasciitilde{}-5°)，保护刀尖，增加强度- 刚度不足：取正值，减小背向力，避免振动
  \item 加工强度高、硬度高材料，同时系统刚性好，选主偏角小- 副偏角：表面质量要求高时取小值
\end{enumerate}
\end{tcolorbox}
\noindent\textbf{次重点与补充：}
\begin{itemize}[leftmargin=2.2em,label=--]
  \item 前角过大：削弱切削刃强度和散热能力，导致刀具磨损加剧，寿命下降。
  \item - 刀具材料：强度韧性好（如高速钢）可取大值。
  \item - 加工性质：粗加工、刚性差取小值；精加工取大值硬质合金外圆车刀前角及后角的推荐值硬质合金外圆车刀前角及后角的推荐值被加工材料 前 角 后角
  \item 钢，铝合金，镁合金，黄铜 25° 8°\textasciitilde{}15°灰铸铁 (<220HBW), 软钢 15° 8°\textasciitilde{}12°碳钢及合金钢(Rm=0.78\textasciitilde{}1.18GPa),可锻铸铁，青铜 10° 8°\textasciitilde{}12°灰铸铁 ( >220HBW), 黄铜(脆性的), 高锰钢 0°\textasciitilde{}5° 8°\textasciitilde{}12°冷硬铸铁，淬硬钢 -10° 8°高硅铸铁 -6°\textasciitilde{}-8° 8°不锈钢 20°\textasciitilde{}25° 10°\textasciitilde{}12°高温合金
  \item 变形 5°\textasciitilde{}8° 14°\textasciitilde{}18°
  \item 铸造 \textasciitilde{}10° \textasciitilde{}10°
  \item 钛合金 5° 10°\textasciitilde{}15°后角的功用及选择
  \item 硬质合金外圆车刀前角及后角的推荐值硬质合金外圆车刀前角及后角的推荐值后角的功用
  \item 增大后角：减小后刀面与工件过渡表面摩擦，使刃口更锋利。
  \item 后角过大：减小切削刃强度和散热能力，易导致刀具损坏。
  \item 。
  \item 斜角切削与刃倾角的选择斜角切削定义
  \item 当刃倾角不为零时，主切削刃与切削速度方向不垂直。此时实际前角增大，显著改善切削条件。
  \item 斜角切削原理速度分解与流屑角示意
  \item 切屑流向影响正负刃倾角下的切屑排出方向对比
  \item 刀刃强度影响不同刃倾角对刀尖保护的作用
  \item 主偏角和副偏角的功用及选择主偏角的功用
  \item - 影响切削层截面形状和切削分力比例（主偏角越小，背向力越大）
  \item - 直接影响刀具寿命和已加工表面粗糙度副偏角的功用
  \item - 减少副切削刃与工件已加工表面的摩擦- 决定残留面积大小，副偏角越小表面粗糙度越低选择原则
  \item - 主偏角：刚度差时取大值，减小背向力防振动；
  \item 粗加工时选较大主偏角，以减小振动，延长刀具寿命；
  \item 主偏角对背向力的影响不同主偏角下切削力的分解情况，
  \item 直观展示背向力的变化趋势。
  \item 副偏角对粗糙度的影响不同副偏角下工件表面残留面积的
  \item 几何形状差异。
  \item 切削用量选择原则
\end{itemize}

\subsection{第五节 切削用量的选择}
\begin{tcolorbox}[colback=yellow!6,colframe=red!55!black,title=重点,left=1mm,right=1mm]
\begin{enumerate}[leftmargin=2.2em,label=\arabic*.]
  \item \textbf{切削用量选择的意义}：提高生产效率、保证生成经济性和切削加工质量
  \item 在条件允许下，先选尽可能大的背吃刀量ap，再根据加工条件和要求选用所运行的最大进给量f ，最后选合理切削速度vc
  \item \textbf{断续切削}：切削表层有硬皮的锻铸件或切削冷硬倾向较为严重的材料（例如不锈钢）时，应尽量使值超过硬皮或冷硬层深
  \item \textbf{进给量 f 的选择方法}：粗加工时，对表面质量没有太高要求，合理的进给量
  \item \textbf{切削速度vc的选择方法}：根据已经选定的背吃刀量、进给量及刀具的寿命，可用公式或查表法确定切削速度。
\end{enumerate}
\end{tcolorbox}
\noindent\textbf{次重点与补充：}
\begin{itemize}[leftmargin=2.2em,label=--]
  \item 切削用量与刀具寿命的关系选择原则
  \item vc ＞f ＞ap切削用量的选择方法
  \item 背吃刀量ap的选择方法背吃刀量根据加工余量确定。
  \item 粗加工时，一次走刀应尽可能切掉全部余量。
  \item 如果是下面几种情况，可几次走刀分切：
  \item 加工余量太大，导致机床动力不足或刀具强度不够工艺系统刚性不足
  \item 度，以防刀具过快磨损。
  \item 应是工艺系统所能承受的最大进给量。
  \item 限制粗加工进给量的因素有：
  \item 机床进给机构的强度刀杆的强度和刚度
  \item 硬质合金或陶瓷刀片的强度等 限制精加工进给量的主要因素是表面粗糙度和加工精度要求。
  \item 实际生产中，经常采用查表法确定进给量。
  \item 精加工时，避开积屑瘤产生的速度区域 断续切削时，要降低切削速度 避开共振频率的速度区域 大件、细长件、薄壁件以及有铸锻外皮工件时，选择较低的切削速度
  \item 车削加工范围
\end{itemize}

\section{第四章 车削加工及普通车床}
\subsection{本章核心知识}
\begin{tcolorbox}[colback=yellow!6,colframe=red!55!black,title=重点,left=1mm,right=1mm]
\begin{enumerate}[leftmargin=2.2em,label=\arabic*.]
  \item 台式钻床：台式钻床结构简单，使用方便，体积小，但只能加工小孔 (一般情况下，孔径<12mm)
  \item 摇臂绕立柱回转并可带着主轴箱沿立柱轴线上下移动，以适应不同高度的工件。主轴箱可通过手动或机动沿摇臂的水平导轨移动。调整主轴 (刀具)的位置，对准所需钻孔的中心而不必移动工件工件可以安装在工作台或底座上。摇臂钻床广泛应用于大型工件或多孔工件的钻削
\end{enumerate}
\end{tcolorbox}
\noindent\textbf{次重点与补充：}
\begin{itemize}[leftmargin=2.2em,label=--]
  \item 立式钻床：工作台可沿立柱上的导轨上下调整位置，以适应不同高度工件的加工立式钻床只适于在单件、小批生产中加工中、小型工件上的孔
\end{itemize}

\subsection{第一节 概述}
\begin{tcolorbox}[colback=yellow!6,colframe=red!55!black,title=重点,left=1mm,right=1mm]
\begin{enumerate}[leftmargin=2.2em,label=\arabic*.]
  \item \textbf{其他加工：如滚花、车槽、切断等。}：应用广泛性：车床应用广泛，占金属切削机床总数的
  \item \textbf{表面成形运动 (核心)}：主运动：工件的旋转运动，提供切削速度，是
  \item \textbf{消耗功率最大的运动。}：进给运动：刀具的移动。车外圆为纵向，车端面为横向，车螺纹为复合螺旋运动。
  \item \textbf{辅助运动 (配合)}：切入运动：使刀具切入工件，达到所需深度。
  \item 总结：车削加工通过工件旋转和刀具移动的相对运动实现。表面成形运动决定工件形状，辅助运动保障加工过程的高效进行。
  \item - 车削加工中心核心特点：
\end{enumerate}
\end{tcolorbox}
\noindent\textbf{次重点与补充：}
\begin{itemize}[leftmargin=2.2em,label=--]
  \item 主要加工回转表面：内外圆柱面、内外圆锥面、成形回转面、回转体的端面。
  \item 孔加工：可以使用钻头、扩孔钻、铰刀等进行钻孔、扩孔和铰孔。
  \item 螺纹加工：车削各种内外螺纹。
  \item 20\%\textasciitilde{}35\%。 图4-1 卧式车床能完成的工作车削加工运动
  \item 快移运动：刀架、尾座的快速移动，用于调整位置。
  \item 空行程运动：如开停机、变速、装卸工件等。
  \item 车床的分类普通车床（非数控）
  \item - 卧式车床及落地车床- 立式车床
  \item - 转塔车床- 仪表车床
  \item - 单轴/多轴自动和半自动车床- 仿形车床和多刀车床
  \item - 专门化车床（如凸轮轴车床、曲轴车床）
  \item 数控车床- CNC车床
  \item 自动化程度高、加工精度高、效率高，正逐步取代传统自动车床。
  \item 随着工业技术的发展，数控车床凭借其高精度和高效率的优势，在现代制造业中正逐渐占据主导地位。
  \item 车床的分类 — 卧式车床车床的分类— 卧式车床
  \item 车床的分类 — 立式车床车床的分类 — 转塔车床
  \item 车床的分类 — 数控车床机床的工艺范围
\end{itemize}

\subsection{第二节 CA6140型卧式车床及传动系统}
\begin{tcolorbox}[colback=yellow!6,colframe=red!55!black,title=重点,left=1mm,right=1mm]
\begin{enumerate}[leftmargin=2.2em,label=\arabic*.]
  \item \textbf{加工范围}：内外圆柱面、圆锥面、环槽、成形面、端面及各种螺纹。还能进行钻孔、扩孔、铰孔、滚花等操
  \item \textbf{加工精度}：- 精车外圆精度：0.01mm - 圆柱度：0.01mm/100mm - 表面粗糙度 Ra：$1.25\sim2.5\,\mu\text{m}$
  \item 总结：CA6140车床结构简单、操作方便，是机械加工中不可或缺的基础设备。
  \item \textbf{进给箱}：安装在床身左端前侧，内部是进给变换机构，用于改变进给量或螺纹导程
  \item \textbf{溜板箱}：位于床身前面，与刀架相连，传递进给运动，实现纵横向进给和快速移动
  \item \textbf{刀架}：由多层滑板组成，用于装夹车刀，并带动刀具实现各种进给运动。
  \item \textbf{尾座}：安装在床身右侧导轨上，可纵向移动，用于支承长工件或安装孔加工刀具
  \item \textbf{件的相对位置精度。}：核心总结：各部件协同工作，共同完成车削加工任务，床身是精度保证的基础。
  \item \textbf{最大工件长度}：可选规格：750/1000/1500/2000mm
  \item \textbf{主轴转速}：正转24级：10\textasciitilde{}1400r/min反转12级：14\textasciitilde{}1580r/min
  \item \textbf{进给量}：纵向64种：0.028\textasciitilde{}6.33mm/r横向64种：0.014\textasciitilde{}3.16mm/r
  \item 车削螺纹范围米制螺纹：44种
  \item 额定功率：7.5kWCA6140车床的传动系统
  \item 系统总览：传动系统图清晰展示了机床的全部运动联系，是理解工作原理的核心。
  \item 主运动传动链：从主电动机传递至主轴，主要实现主轴的旋转运动及变速换向功能。
  \item 核心原理：通过改变传动链中齿轮的啮合关系，实现主轴不同转速和刀架不同进给量。
  \item \textbf{主运动传动链}：图4-3 CA6140型卧式车床传动系统图
  \item \textbf{传动路线}：电动机 -> V带轮 -> 主轴箱（轴I-VI）-> 主轴。
  \item \textbf{换向机构}：轴I上的双向片式摩擦离合器M1，精准控制主轴的正转、反转和停止状态。
  \item \textbf{变速机构}：通过轴I-II、II-III、III-IV-V等多组滑移齿轮的不同啮合组合，实现多级变速。
  \item \textbf{转速级数}：正转24级（6高18低）；反转12级（高转速退
  \item \textbf{公用段：主轴 -> 进给箱}：分支路径：进给箱后分为丝杠与光杠两路
  \item \textbf{分类与功能特点}：丝杠传动：实现车削螺纹（高精度传动）
\end{enumerate}
\end{tcolorbox}
\noindent\textbf{次重点与补充：}
\begin{itemize}[leftmargin=2.2em,label=--]
  \item 作。
  \item 应用场景通用性强，但自动化程度较低。
  \item 广泛适用于单件、小批量生产以及各类维修车间。
  \item 机床的布局与组成主轴箱
  \item 安装在床身左端，内部有主轴和变速机构，实现主运动的变速和换向。
  \item 。
  \item 床身机床的基础部件，支承并保证其他部
  \item CA6140车床的主要技术性能最大回转直径
  \item 床身上：400mm刀架上：210mm
  \item 寸制螺纹：20种主电动机功率
  \item 进给运动传动链：从主轴传递至刀架，驱动刀架完成进给运动和车螺纹运动。
  \item 动力逐级传递，确保高效运行。
  \item 刀，节省时间）。
  \item 进给运动传动链传动链组成与路径
  \item 两末端件：主轴（起点）-> 刀架（终点）
  \item 光杠传动：实现纵向/横向机动进给机动进给：包含纵向（前后）和横向（左右）
  \item 进给运动传动链组成框图其他常用车床—立式车床
\end{itemize}

\subsection{第四节 其他常用车床与刀具}
\begin{tcolorbox}[colback=yellow!6,colframe=red!55!black,title=重点,left=1mm,right=1mm]
\begin{enumerate}[leftmargin=2.2em,label=\arabic*.]
  \item 加工过程中，多工位刀架可周期性地转位，将不同刀具依次转到加工位置，对工件进行加工转塔车床的优点是:在成批生产中，特别是在加工形状复杂的工件时，生产效率高车床刀具—种类和用途
  \item \textbf{车床刀具—常用结构}：整体车刀：刀杆与刀头为一整体，高速钢材
  \item \textbf{料为主，可重新刃磨}：焊接车刀：在碳钢刀杆（常用 45 钢）上根据刀片的形状和尺寸铣出刀槽后，将硬质合金刀片钎焊在刀槽中，可重新刃磨焊接装配车刀：将硬质合金刀片钎焊在小刀块上，再将小刀块装配到刀杆上，多用于重
  \item \textbf{型车刀。}：机夹重磨式车刀：夹重磨式车刀的刀片和刀杆是两个可拆卸的独立元件，工作时靠夹紧装置把它们固定在一起，可重新刃磨焊机夹可转位式车刀：将压制成一定几何参数的多边形刀片（如硬质合金刀片），用机
  \item \textbf{成形车刀}：加工回转体成形表面的专用刀具，其刃形是
  \item 根据零件的廓形设计的具有以下特点：
  \item \textbf{生产率高，一般用于大批量生产。}：加工质量和精度稳定，加工零件精度可达 IT8～IT7，表面粗糙度 Ra 可达 $10\sim2.5\,\mu\text{m}$。
  \item \textbf{刀具使用寿命长}：成形车刀磨损后，一般只
  \item \textbf{需重磨前刀面，刃磨方便。}：设计、制造复杂，成本较高 成形车刀不宜
\end{enumerate}
\end{tcolorbox}
\noindent\textbf{次重点与补充：}
\begin{itemize}[leftmargin=2.2em,label=--]
  \item 轴向尺寸大而径向尺寸比较小的大型和超重型零件，如各类盘、轮类零件其他常用车床—转塔车床
  \item 械夹固的方法装夹在标准的刀杆上。当刀片
  \item 上一个切削刃用钝后，松开夹紧机构，将刀片转位换成另一个新的切削刃车床刀具—成形车刀
  \item 用于单件小批量生产铣削加工工艺范围
\end{itemize}

\subsection{第二节 钻削加工及设备}
\begin{tcolorbox}[colback=yellow!6,colframe=red!55!black,title=重点,left=1mm,right=1mm]
\begin{enumerate}[leftmargin=2.2em,label=\arabic*.]
  \item 钻床的主要类型镗削加工范围
\end{enumerate}
\end{tcolorbox}

\section{第五章 其他典型切削加工方法与设备}
\subsection{本章核心知识}
\textit{（本节暂无可提取文本）}

\subsection{第一节 铣削加工与设备}
\begin{tcolorbox}[colback=yellow!6,colframe=red!55!black,title=重点,left=1mm,right=1mm]
\begin{enumerate}[leftmargin=2.2em,label=\arabic*.]
  \item \textbf{加工范围广泛}：可加工平面、沟槽、螺旋形表面、齿轮、回转体表面及内孔，也能进行切断作业。
  \item \textbf{加工精度较高}：精铣表面粗糙度 Ra 可达 $3.2\sim1.6\,\mu\text{m}$，尺寸精度可
  \item \textbf{切削效率高}：采用多刃刀具切削，同时参与切削的刀刃多，生
  \item \textbf{端铣法}：使用面铣刀(刀齿分布在铣刀端面上)
  \item 逆铣特点：铣刀旋转切入方向与工件进给方向相反。
  \item 顺铣特点：铣刀旋转切入方向与工件进给方向相同。
  \item 注意：需消除进给机构间隙，防止工作台窜动。
  \item \textbf{不对称逆铣}：切入厚度小，切出厚度大，可减少刀齿冲击，
  \item \textbf{不对称顺铣}：切入厚度大，切出厚度小，适用于加工不锈钢
  \item \textbf{等材料，减小刀具破损。}：图5-4 对称铣削与不对称铣削示意图
  \item \textbf{应用要点}：根据材料特性和加工要求选择切入方式，可优化切
  \item 铣刀的线速度，计算公式为：v = pi d n / 1000d：铣刀直径 n：铣刀转速进给量 ( f )
  \item \textbf{立式升降台铣床}：主轴垂直布置，主要用于加工平面、沟槽、台阶等，操作方便，通用性强。
  \item \textbf{大型高效设备，具有龙门式框架。}：用于加工大型工件的平面和沟槽，可多刀
\end{enumerate}
\end{tcolorbox}
\noindent\textbf{次重点与补充：}
\begin{itemize}[leftmargin=2.2em,label=--]
  \item 达IT7\textasciitilde{}IT9级标准。
  \item 产率显著高于单刃刀具。
  \item 铣削方式 - 周铣与端铣周铣法
  \item 使用圆柱铣刀(刀齿分布在铣刀圆柱面上)进行加工，铣刀轴线与被加工表面平行，主要用于加工狭长平面。
  \item 进行加工，铣刀轴线与被加工表面垂直，主要用于加工较大的平面。
  \item 周铣和端铣加工示意图对比铣削方式 - 逆铣与顺铣
  \item 效果：切入时切削厚度为零，表面易产生划擦，刀具磨损较快。
  \item 适用：工作平稳，适合加工带硬皮的工件。
  \item 效果：切入时切削厚度最大，无划擦，刀具寿命长，表面质量好。
  \item 逆铣(a)与顺铣(b)切削原理对比应用建议
  \item 在实际加工中，需根据工件状态（如是否有硬皮）、机床刚性及表面质量要求灵活选择铣削方式。现代数控机床通常具备消除间隙的功能，更适合采用顺铣以提高效率和质量。
  \item 铣削方式 - 对称与不对称铣削对称铣削
  \item 铣刀相对于工件的位置对称，刀齿切入和切出的切削厚度相同。
  \item 使切削更平稳。
  \item 削过程，显著提高加工质量和刀具寿命。
  \item 铣削四要素铣削速度 (v)
  \item - 每齿进给量( fz )：每转过一个刀齿的位移 mm/z- 每转进给量( f )：每转一转的位移 mm/r- 进给速度( vf )：单位时间进给量 mm/minvf = f n = fz z n铣削深度 (ap)
  \item 平行于铣刀轴线测量的切削层尺寸。
  \item 铣削宽度 (ae)
  \item 垂直于铣刀轴线测量的切削层尺寸。
  \item 铣刀的类型及用途加工平面的铣刀
  \item 方肩铣：同时生成相互垂直的两个面，这需要将周铣与端铣相结合
  \item 周铣：使用圆柱铣刀(刀齿分布在铣刀圆柱面上)进行加工
  \item 端铣：使用面铣刀(刀齿分布在铣刀端面上)进行加工
  \item 加工沟槽的铣刀加工圆弧面的铣刀
  \item 铣床的主要类型卧式升降台铣床
  \item 主轴水平布置，应用最广泛的铣床类型适用于单件和小批量生产
  \item 卧式升降台铣床 立式升降台铣床床身式铣床
  \item 工作台不升降，依靠铣头移动完成进给。
  \item 刚度高，生产效率高，适合成批生产龙门铣床
  \item 同时切削钻削加工工艺范围
\end{itemize}

\subsection{第二节 钻削加工及设备}
\begin{tcolorbox}[colback=yellow!6,colframe=red!55!black,title=重点,left=1mm,right=1mm]
\begin{enumerate}[leftmargin=2.2em,label=\arabic*.]
  \item \textbf{加工特点：}：- 钻头是定尺寸刀具，加工精度相对较低- 由于钻头的刚性和导向性较差，加工时容易产生引偏，导致孔的轴线歪斜
  \item 麻花钻是最常用的孔加工刀具，一般用于实体材料上孔的粗加工；钻孔的尺寸精度为 IT11～IT13，表面粗糙度 Ra 为 $50\sim12.5\,\mu\text{m}$。麻花钻的切削部分和导向部分：
  \item \textbf{钻孔刀具 – 扩孔钻}：扩孔钻是用来对工件上已有孔进行扩大加工的刀具。
  \item 扩孔后，孔的精度可达到 IT9\textasciitilde{}IT10，表面粗糙度 Ra 可达到 $6.3\sim3.2\,\mu\text{m}$。结构特点：
  \item \textbf{钻孔刀具 – 铰刀}：铰刀是一种半精加工或精加工孔的常用刀具；铰孔后孔的精度可达 IT7\textasciitilde{}IT9，表面粗糙度 Ra 可达 $1.6\sim0.4\,\mu\text{m}$
  \item \textbf{铰孔工艺及其应用}：铰孔余量对铰孔质量影响很大，余量太大，铰刀负荷大；余量下太小，不能达到精
  \item \textbf{加工的目的；}：铰孔采用较低的切削速度，避免积屑瘤必须使用适当的切削液进行冷却、润滑和清洗与镗孔和磨孔相比，生产效率高
\end{enumerate}
\end{tcolorbox}
\noindent\textbf{次重点与补充：}
\begin{itemize}[leftmargin=2.2em,label=--]
  \item - 表面粗糙度值较大钻孔刀具 – 麻花钻
  \item 切削部分由两条主切削刃、两条副切削刃、一条横刃及两个前刀面和两个后刀面组成导向部分有两条螺旋槽和两条棱边，螺旋槽起排屑和输送切削液的作用，棱边起导向、修光孔的作用。
  \item 导向部分有微小的倒锥度，从前端到尾部每 100mm 长度上直径减少 0.03～0.12mm，以减少与孔壁的摩擦
  \item 扩孔钻与麻花钻相似，但刀齿数较多，没有横刃扩孔钻是用来对工件上已有孔进行扩大加工的刀具。
  \item 扩孔后，孔的精度可达到 IT9\textasciitilde{}IT10，表面粗糙度 Ra 可达到 $6.3\sim3.2\,\mu\text{m}$。扩孔除了可以加工圆柱孔之外，还可以用各种特殊形状的扩孔钻（亦称锪钻）来加工各种沉头座孔和锪平端面
  \item 铰刀的刀齿数多 (4\textasciitilde{}12个齿)，加工余量小，导向性好，刚度大。
  \item 钻床的主要类型
\end{itemize}

\subsection{第三节 镗削加工及设备}
\begin{tcolorbox}[colback=yellow!6,colframe=red!55!black,title=重点,left=1mm,right=1mm]
\begin{enumerate}[leftmargin=2.2em,label=\arabic*.]
  \item \textbf{加工特点：}：镗削主要用于工件上已有铸造孔或加工孔的后续加工，常用于加工尺寸较大及精度要求较高的场合，特别适宜加工尺寸精度和位置精度要求十分严格的孔系。镗削是常用的加工方法，其加工范围很广，既可进行粗加工，也可进行精加工。镗削加工的标准公差等级一般为 IT6\textasciitilde{}IT7，表面粗糙度 Ra 一般为 $6.3\sim0.8\,\mu\text{m}$
  \item 镗刀 – 单刃镗刀结构特点：
  \item 单刃镗刀的刀头结构与车刀类似，使用时用紧固螺钉将其装夹在镗杆上通孔单刃镗刀
  \item \textbf{阶梯孔单刃镗刀 盲孔单刃镗刀}：双刃镗刀两端都有切削刃，工作时基本上可消除径向力对镗杆的影响
  \item \textbf{加工范围：}：卧式镗床加工范围广泛，除镗孔外，还可以车端面、铣平面、车外圆、车螺纹及钻孔等
  \item \textbf{镗孔方式：}：工件旋转，镗刀直线进给。适合加工回转体类零件上的孔(与外圆表面有同轴度要求)
  \item 典型镗床 – 坐标镗床定义及加工范围：
\end{enumerate}
\end{tcolorbox}
\noindent\textbf{次重点与补充：}
\begin{itemize}[leftmargin=2.2em,label=--]
  \item 双刃镗刀大多采用浮动结构典型镗床 – 卧式镗床
  \item 镗刀做旋转运动（主运动）， 工件做直线进给运动；
  \item 适合加工箱体等非回转类零件，或者大型工件上的孔(与孔所在的端面有垂直度等要求)
  \item 刀具既旋转又进给的镗孔方式适合加工较短的孔
  \item 坐标镗床具有坐标位置精密测量装置，能精确地确定工作台、主轴箱等移动部件的位移量，实现工件和刀具的精确定位用于加工孔本身精度及位置精度要求都很高
  \item 的孔系，如钻模、镗模和量具等零件上的精密孔
  \item 坐标镗床的工艺范围很广，除镗孔、钻孔、扩孔、铰孔及精铣平面沟槽外，还可进行精密刻线和划线，以及进行孔距和直线尺寸的精密测量工作
  \item 立式单柱坐标镗床立式双柱坐标镗床
  \item 齿轮加工概述
\end{itemize}

\subsection{第四节 齿轮加工及设备}
\begin{tcolorbox}[colback=yellow!6,colframe=red!55!black,title=重点,left=1mm,right=1mm]
\begin{enumerate}[leftmargin=2.2em,label=\arabic*.]
  \item 定义：加工齿轮齿面的方法，是机械制造中获得高精度齿轮的关键工艺环节。
  \item 原理：使用与齿轮齿槽形状相符的成形刀具进行加工（如铣齿）。
  \item 特点：工艺简单，成本较低，但加工精度相对较低，适用于单件小批量生产。
  \item 原理：利用齿轮啮合原理，通过刀具与工件的啮合运动加工（如滚齿、插齿、磨齿）。
  \item 特点：加工精度高，效率高，是现代齿轮加工的主要方法，适用于大批量生产。
  \item 总结：从精度和效率综合考量，展成法因其高适应性和高精度，已成为现代工业齿轮加工的主流技术。
  \item 齿轮加工原理成形法：(铣齿、插削、磨削)
  \item \textbf{而形成的。}：在展成法加工齿轮时，用同一把刀具可以加工相同模数的任意齿数的齿轮。
  \item \textbf{轴向进给运动}：主运动 即插齿刀的往复直线运动分齿运动(展成运动) 即维持插齿刀与被切齿轮
\end{enumerate}
\end{tcolorbox}
\noindent\textbf{次重点与补充：}
\begin{itemize}[leftmargin=2.2em,label=--]
  \item 成形法 (Forming Method)
  \item 展成法 (Generating Method)
  \item 成形法加工齿轮是使用切削刃形状与被切齿轮的齿槽形状完全相符的成形刀具切出齿轮的方法
  \item 一般用于铣床上用盘形铣刀或指形齿轮铣刀铣削齿轮，也可以在刨床或插床上用成形刀具刨削、插削齿轮
  \item 展成法展成法加工齿轮是利用齿轮啮合的原理进行
  \item 的，把其中一个齿轮转化为刀具，另一个转化为工件，并强制刀具和工件做严格的啮合运动，被加工工件的齿形表面是在刀具和工件包络过程中由刀具切削刃的位置连续变化
  \item 其加工精度和生产率都比较高，在齿轮加工中应用最为广泛
  \item 滚齿刀插齿刀
  \item 剃齿刀展成法 – 滚齿
  \item 原理运动
  \item 滚齿刀的旋转运动 主运动工件的旋转运动 分齿运动
  \item 之间啮合关系的运动圆周进给运动
  \item 径向进给运动让刀运动
  \item 展成法 – 剃齿剃齿刀的正反转动
  \item 工作台轴向进给运动展成法 – 磨齿
  \item 砂轮滚齿、分度磨齿齿轮精度6级
  \item 电火花加工（EDM）
\end{itemize}

\subsection{第五节 特种加工}
\begin{tcolorbox}[colback=yellow!6,colframe=red!55!black,title=重点,left=1mm,right=1mm]
\begin{enumerate}[leftmargin=2.2em,label=\arabic*.]
  \item \textbf{加工原理及特点}：电火花加工是利用工具电极和工件电极之间的脉冲性电火花放电产生的高温去除工件上多余的材料，使工件获得预定的尺寸和表面
  \item 电火花加工类型：
  \item \textbf{超声加工}：超声波加工是利用工具端面的超声频振动（振动频率为19000～25000Hz )，驱动工作
  \item \textbf{纳米制造}：纳米加工技术是以纳米级精度对材料进行原子或分子尺度去除、搬迁及重组的技术“自上而下”的高精度加工（如纳米光刻、纳
  \item \textbf{米压印）}：“自下而上”的直接构建（原子层沉积）
\end{enumerate}
\end{tcolorbox}
\noindent\textbf{次重点与补充：}
\begin{itemize}[leftmargin=2.2em,label=--]
  \item 粗糙度要求。
  \item 工件导电电火花加工（EDM）
  \item 电火花穿孔加工电火花成型加工
  \item 电火花线切割加工电化学加工（ECM）
  \item 利用金属在电解液中受到电化学阳极溶解，将工件加工成形
  \item 液中的悬浮磨料撞击加工表面的加工方法。
  \item 超声加工不仅能加工硬质合金、淬火钢等脆硬金属材料，而且更适合于加工玻璃、陶瓷、半导体锗和硅片等不导电的非金属脆硬材料，同时还可以用于清洗、焊接和探伤等磁流变抛光技术
  \item 在外加磁场的作用下，磁流变液体的黏度会随着磁场的增强而增大，形成类固体的结构从而具有较高的屈服强度。在磁流变液中加入磨料，利用磁流变液的固化作用来对工件表面进行光整加工
  \item 加工大口径平面、球面、离轴非球面等光学元件
  \item 砂轮的特性与选择 – 普通砂轮的特性
\end{itemize}

\section{第五章 磨 削}
\subsection{本章核心知识}
\begin{tcolorbox}[colback=yellow!6,colframe=red!55!black,title=重点,left=1mm,right=1mm]
\begin{enumerate}[leftmargin=2.2em,label=\arabic*.]
  \item \textbf{砂轮磨削区温度分布示意}：磨粒磨削点温度 (thetadot)
  \item 磨粒切削刃与切屑接触部分的温度，是磨削中温度最高的部位，也是磨削热的热源，可达1000°C以上。不但影响工件表面质量，且与磨粒的磨损及切屑熔着密切相关。
  \item \textbf{砂轮磨削区温度 (thetaA)}：砂轮与工件接触区域的平均温度。与与磨削烧伤、磨削裂纹的产生有密切关系
  \item \textbf{工件平均温度}：指磨削热传入工件而引起的工件温升,其影响工件的形状和尺寸精度。
  \item \textbf{淬火烧伤 超过Ac3}：如果此时冷却充分，则表层将急冷形成二次淬火马氏体组织。工件表层硬度较原来的回火马氏体高，但很薄，其下层因冷却速度慢仍为硬度较低的回火索氏体和屈氏体
  \item \textbf{回火烧伤 未超过Ac3}：但超过马氏体的转变温度，这时工件表层的组织将转变成硬度较低的回火屈氏体或索
  \item \textbf{冷却方法}：加大切削液流量和采用喷雾冷却、高压冷却等加速热量的传出，降低磨削区温度，可
\end{enumerate}
\end{tcolorbox}
\noindent\textbf{次重点与补充：}
\begin{itemize}[leftmargin=2.2em,label=--]
  \item 在精密磨削时，要尽可能降低工件的平均温度并防止局部温度不匀
  \item 退火烧伤 超过Ac3马氏体转变为奥氏体此时无切削液，则表面层硬度急剧下降，工件表面层被退火，故这种烧伤称为退火烧伤。工件干磨时常发生这种情况。
  \item 氏体磨削用量
  \item 当v和 fr较小时，不易出现烧伤。
  \item 以有效避免烧伤现象砂轮接触长度
  \item 砂轮与工件的接触长度大，易堵塞而不易冷却，容易出现烧伤
\end{itemize}

\subsection{第一节 砂轮的特性与选择}
\begin{tcolorbox}[colback=yellow!6,colframe=red!55!black,title=重点,left=1mm,right=1mm]
\begin{enumerate}[leftmargin=2.2em,label=\arabic*.]
  \item \textbf{磨料 (Abrasive)}：砂轮的主要组成成分，必须具备高硬度，决定砂轮的切削能力与耐用度
  \item \textbf{棕刚玉(A) / 白刚玉(WA)}：韧性较好，主要用于磨削碳钢、合金钢等
  \item \textbf{碳化硅类 (SiC)}：黑碳化硅(C) / 绿碳化硅(GC)
  \item \textbf{粗粒度特点}：粒度数值小，磨削效率高，但表面粗糙度大，适用于粗磨
  \item \textbf{细粒度特点}：粒度数值大，磨削效率低，但表面光洁度高，适用于精磨
  \item \textbf{选择原则}：磨软金属选粗粒度，磨硬脆金属选细粒度。
  \item 陶瓷结合剂 (V)：强度高、耐热性好，应用最广。
  \item \textbf{用于切断和开槽}：根据磨床类型和具体加工要求选择合适的砂轮形状与组织
  \item 核心定义：当砂轮线速度超过150m/s时，称为超高速磨削。
  \item 结构特点：采用高强度基体（如碳纤维复合材料），仅边缘附着薄层磨料，以承受巨大离心力。
  \item 关键技术：配备在线自动平衡系统，实时补偿不平衡量，保证稳定性与安全性。
  \item 核心优势：相比传统工艺，大幅提高磨削效率和加工表面质量。
\end{enumerate}
\end{tcolorbox}
\noindent\textbf{次重点与补充：}
\begin{itemize}[leftmargin=2.2em,label=--]
  \item 刚玉类 (Al2O3)
  \item 钢材，应用最广泛。
  \item 硬度更高、性脆锋利，适用于磨削铸铁、硬质合金、有色金属等。
  \item 高硬磨料CBN / 人造金刚石(MD)
  \item 超硬材料，用于加工高硬度、难加工材料，如陶瓷、高速钢、钛合金。
  \item 不同磨料的选择取决于工件材料的硬度、韧性及加工要求砂轮的特性与选择 – 普通砂轮的特性
  \item 粒度 (Grit) 决定加工表面粗糙度与效率粒度定义
  \item 指磨料颗粒的大小，分为磨粒(F4-F220)和微粉(F230-F1200)。
  \item 硬度 (Hardness)
  \item 指磨粒从砂轮表面脱落的难易程度结合剂 (Bond)
  \item 粘结磨粒，决定砂轮强度与韧性砂轮硬度并非指磨料本身，而是指磨粒在受力下从砂轮表面脱落的难易程度。
  \item 加工硬材料 -> 选软砂轮磨粒易脱落，保持砂轮锋利度
  \item 加工软材料 -> 选硬砂轮磨粒不易脱落，提高耐用度
  \item 结合剂是将松散的磨粒粘结在一起，形成具有一定形状和强度砂轮的材料。
  \item 树脂结合剂 (B)：弹性好，多用于薄片砂轮。
  \item 橡胶结合剂 (R)：强度高弹性大，用于切断砂轮。
  \item 总结：硬度决定自锐性，结合剂决定性能。合理选择是高效加工的关键。
  \item 组织 (Structure)
  \item 磨料、结合剂与空隙的体积比例形状尺寸 (Shape)
  \item 根据机床与加工对象需求选择紧密组织：磨粒占比高，空隙小，适合精磨加工疏松组织：空隙占比大，排屑散热好，适合粗磨及软材平形砂轮
  \item 用于外圆、平面磨削杯形砂轮
  \item 用于磨削平面、刀具碗形砂轮
  \item 用于刃磨刀具薄片砂轮
  \item 砂轮的特性与选择 – 超高速磨削砂轮
  \item 超高速砂轮结构与实物示例磨削过程
\end{itemize}

\subsection{第二节 磨削运动与磨削过程}
\begin{tcolorbox}[colback=yellow!6,colframe=red!55!black,title=重点,left=1mm,right=1mm]
\begin{enumerate}[leftmargin=2.2em,label=\arabic*.]
  \item \textbf{01. 滑擦阶段}：磨粒刚接触工件，仅产生弹性和塑性变形，无切屑产生，主要发生摩擦生热
  \item \textbf{02. 刻划阶段}：磨粒切入工件，材料被挤压向两侧隆起形成划痕，仍未形成完整切屑
  \item \textbf{03. 切削阶段}：切削厚度足够大，材料被剪切形成磨屑，并沿磨粒前刀面顺利流出核心机理：由于砂轮表面磨粒高度参差不齐，实际磨削加工中，这三个阶段是同时存在且动态变化的。
  \item \textbf{磨削运动}：主运动 (v) - 砂轮高速旋转提供主要切削速度，是磨削加工的核心动力来源径向进给运动 (fr) - 砂轮切入砂轮向工件方向的切入，决定单次磨削的深度（背吃刀量）
  \item \textbf{1000×60}：v——磨削速度 (m/s);d0——砂轮直径 (mm);n0——砂轮转速(r/ min)。
\end{enumerate}
\end{tcolorbox}
\noindent\textbf{次重点与补充：}
\begin{itemize}[leftmargin=2.2em,label=--]
  \item 轴向进给运动 (fa) - 相对移动工件相对于砂轮的轴向移动，决定磨削的宽度范围工件运动 (vw) - 旋转/移动工件自身的旋转或直线移动，提供进给速度以配合切削$v= \pi d_0 n_0$
  \item 磨削力与磨削功率
\end{itemize}

\subsection{第三节 磨削力、磨削功及磨削温度}
\begin{tcolorbox}[colback=yellow!6,colframe=red!55!black,title=重点,left=1mm,right=1mm]
\begin{enumerate}[leftmargin=2.2em,label=\arabic*.]
  \item - 径向力 (Fp)：最大，通常是切向力 Fc 的1.6-3.2倍，是导致工艺系统变形的主要原因。
  \item - 切向力 (Fc)：决定磨削功率的关键因素，用于计算功率消耗。
  \item - 轴向力 (Ff)：通常较小，对加工影响有限，一般可忽略不计。
  \item 由切向力 (Fz) 和砂轮线速度 (v) 共同决定，计算公式如下：
\end{enumerate}
\end{tcolorbox}
\noindent\textbf{次重点与补充：}
\begin{itemize}[leftmargin=2.2em,label=--]
  \item 磨削力分解 (Grinding Force)
  \item 磨削功率计算 (Grinding Power)
  \item Pm = Fz ×v / 1000(kW)
  \item 磨削温度磨削温度对工件表面质量的影响
  \item 影响磨削烧伤的工艺因素高精度低表面粗糙度磨削
\end{itemize}

\subsection{第四节 先进的磨削方法}
\begin{tcolorbox}[colback=yellow!6,colframe=red!55!black,title=重点,left=1mm,right=1mm]
\begin{enumerate}[leftmargin=2.2em,label=\arabic*.]
  \item 核心工艺技术要点精细修整砂轮：磨粒等高微刃
  \item \textbf{小进给量：}：0.005\textasciitilde{}0.025mm/r低磨削速度：15\textasciitilde{}30m/s
  \item \textbf{实现高精度与高效率的完美统一}：高效磨削 - 缓进给磨削（深磨工艺）
  \item \textbf{核心定义：大切深 + 缓进给}：- 切深可达30mm（一次成形） | 进给5\textasciitilde{}300mm/min
  \item \textbf{工艺优势：高效且全能}：- 效率：以磨代车/铣，毛坯直达成品- 寿命：低冲击，砂轮耐用 | 适合难加工材料
  \item \textbf{挑战与对策：防烧伤}：- 风险：接触弧长易烧伤 | 对策：大气孔砂轮+高压冷却
  \item - 点接触：砂轮轴线倾斜，实现理论点接触- 超硬砂轮：CBN/金刚石材质，厚度仅4\textasciitilde{}6mm- 双高速叠加：实际磨削速度可达250m/s工艺核心优势
  \item - 力小温低：适合细长轴，实现“冷态”加工- 高效柔性：效率提升6倍，一次装夹多工序加工原理示意
  \item \textbf{工艺分类}：- 闭式：高速回转，效率极高，适合粗加工- 开式：低速运行，冷却充分，表面质量优
  \item \textbf{砂带结构}：- 基体：布、纸或化纤织物，承载磨料- 磨料：静电植砂工艺，确保磨粒等高性
  \item \textbf{技术优势}：加工效率高、散热性能好、工件变形小、适应性极强
  \item \textbf{广泛的工业应用}：涵盖金属加工、木材加工及航空航天等高端制造领域，是提升生产效率的关键工艺。
\end{enumerate}
\end{tcolorbox}
\noindent\textbf{次重点与补充：}
\begin{itemize}[leftmargin=2.2em,label=--]
  \item 极致表面质量：粗糙度 Ra $<0.4\,\mu\text{m}$
  \item 高效加工替代传统研磨工艺
  \item 光磨行程：后期多次无进给严苛的设备要求
  \item 高刚度高精度磨床，工作台无爬行，具备精确微量进给机构加工效果指标
  \item Ra $<0.4\,\mu\text{m}$（镜面级）
  \item 高效磨削 - 高速、超高速磨削高速磨削
  \item 线速度 > 45m/s超高速磨削
  \item 线速度 > 150m/s生产率极高
  \item 磨除率提升30\%\textasciitilde{}100\%砂轮寿命长
  \item 单磨粒负荷小，损耗低表面质量优
  \item 切削厚度薄，粗糙度低磨削力更小
  \item 径向力减小，工件不变形磨削温度低
  \item 热损伤风险大幅降低超硬磨料适配
  \item 完美发挥CBN/金刚石性能技术核心价值
  \item 高效磨削 – 快速点磨核心技术特征
  \item 图示展示了砂轮倾斜安装状态，以及其与工件微小的点接触区域，这是实现高速高效磨削的关键几何条件。
  \item 砂带磨削核心定义
  \item 以砂带代替传统砂轮的高效磨削方式，通过高速回转实现材料加工。
  \item 外圆磨床
\end{itemize}

\subsection{第五节 磨 床}
\begin{tcolorbox}[colback=yellow!6,colframe=red!55!black,title=重点,left=1mm,right=1mm]
\begin{enumerate}[leftmargin=2.2em,label=\arabic*.]
  \item \textbf{核心用途}：主要用于磨削轴类零件的外圆柱面和圆锥面，是机械加工车间实现高精度回转面加工的关键。
  \item \textbf{两种核心磨削工艺}：- 纵磨法：工件往复移动，精度高，适合精磨。
  \item \textbf{万能外圆磨床特性}：工作台可偏转，砂轮架可回转，可扩展加工内孔
  \item \textbf{应用价值总结}：兼顾高精度与高柔性，一机多用特性使其成为机械加工车间不
  \item 卧轴矩台 / 立轴矩台 / 立轴圆台 / 卧轴圆台核心价值：高精度 · 高光洁度 · 强通用性无心外圆磨床
  \item \textbf{核心工作原理}：工件置于砂轮与导轮之间由托板支承，导轮直接控制工件的旋转速度与进给运动，无需顶尖装夹。
  \item \textbf{显著技术优势}：装夹便捷可连续磨削，生产效率高；消除定位误差，
  \item \textbf{两种磨削方法}：纵磨法（自动通过磨削区，适合无台阶）；
\end{enumerate}
\end{tcolorbox}
\noindent\textbf{次重点与补充：}
\begin{itemize}[leftmargin=2.2em,label=--]
  \item - 切入磨法：工件不移动，效率高，适合批量粗磨。
  \item 及复杂台阶面。
  \item 可或缺的关键设备。
  \item 外圆磨床平面磨床
  \item 专注各类工件平面的高精度磨削，确保加工表面的平整度与光洁度。
  \item 两种主流磨削方式对比周边磨削
  \item 精度高 | 效率低端面磨削
  \item 效率高 | 精度略低设备常见分类
  \item 加工精度高，一致性好。
  \item 横磨法（横向进给，适合带台阶）。
  \item 机械制造技术基础C 第八章 机械加工工艺规程的制订
\end{itemize}

\subsection{第一节 基本概念}
\begin{tcolorbox}[colback=yellow!6,colframe=red!55!black,title=重点,left=1mm,right=1mm]
\begin{enumerate}[leftmargin=2.2em,label=\arabic*.]
  \item 生产过程 (Production)
  \item 生产过程是将原材料或半成品转变为成品的全部过程。
  \item 它是产品决策 、设计、毛坯制造 (在锻压车间或铸造车间进行 )、零件的机械加工 、零件的热处理、机械装配、产品调试、检验和试车、包装等一系列相互关联的劳动过程的总和
  \item 工艺过程：改变生产对象的形状、尺寸、相互位置和性质，使其成为成品或半成品的过程。
  \item \textbf{机械加工工艺过程 (核心)}：用切削方法改变毛坯形状、尺寸和性能，制造合格零件。
  \item 多工位加工可提升生产率阶梯轰加工工艺过程示例：工序-安装-工位-工步
  \item 单件、小批生产首选定义：通过工具和仪表找出工件正确位置的过程。
  \item \textbf{分类：}：- 直接找正：用千分表、划针等直接找正- 划线找正：先划线，再按线找正
  \item 定义：利用夹具对工件进行定位、夹紧的方法。
  \item \textbf{核心优势：}：- 操作方便、迅速，精度高且稳定- 单件小批生产可使用组合夹具
  \item 特点：单件小批生产首选调整法
  \item 特点：成批/大量生产定尺寸刀具法
  \item 特点：依赖设备精度生产类型 – 生产纲领
  \item 生产纲领定义：生产纲领是指企业在计划期内应当生产的产品产量
  \item \textbf{零件在计划期为一年的生产纲领：}：N=Qn(1+a\%)(1+b\%)
  \item 单件生产特点：品种多、产量少，不定期重复。
  \item 成批生产特点：分批生产，周期性交替重复。
  \item 分类：小批、中批、大批生产。
  \item 大量生产特点：产量大，工序高度重复，自动化。
  \item \textbf{毛坯的选择}：根据零件的工作情况选择铸件、锻件、焊件、棒料等
  \item \textbf{工序设计}：加工简图的绘制，加工余量、工序尺寸和公差的计
  \item \textbf{算，技术要求的确定，机床设备、切削用量、时间}：定额和工艺装备的选择确定是工序设计的内容
  \item 概念：零件结构工艺性是指所设计的零件在满足使用要求的前提下制造的可行性和经济性准备阶段
  \item \textbf{工艺过程拟订}：确定生产类型及毛坯的种类 划分加工阶段，安排工序的集中和分散确定 各个工序的机床设备及工艺装备
  \item \textbf{填写工艺文件}：机械加工零件结构的工艺性分析合理标注零件的尺寸、公差和表面粗糙度值：
  \item \textbf{以形状简单的轮廓要素为基准标注尺寸，避免尺寸换算。}：零件上的尺寸公差、几何公差和表面粗糙度值的标注，应根据零件的功能经济、合理地决定
  \item \textbf{零件形状的工艺性：}：形状尽量简单，规格尽量标准和统一 尽量采用普通设备和标准刀具加工 加工面与非加工面应分开
\end{enumerate}
\end{tcolorbox}
\noindent\textbf{次重点与补充：}
\begin{itemize}[leftmargin=2.2em,label=--]
  \item 机械装配工艺过程(零件组装成部件/机器)
  \item 生产过程和机械加工工艺过程机械制造技术基础C 第八章 机械加工工艺规程的制订
  \item 机械加工工艺过程的组成工序 (Process)
  \item 工序是指一个或一组工人，在一台机床或一个工作地对一个或同时对几个工件所连续完成的那一部分工艺过程。
  \item 生产计划/成本核算的基本单元工步 (Step)
  \item 加工表面和工具不变时，连续完成的那一部分工序。
  \item 工序的细分组成单位安装 (Setup)
  \item 工件经一次装夹后所完成的那一部分工序。
  \item 应尽可能减少安装次数工位 (Station)
  \item 一次装夹后，工件与夹具相对刀具占据的每个位置。
  \item 工序 01：车削安装：1次 (自定心卡盘)
  \item 工位：1个工步详情：
  \item 1.车端面A2.车外圆F
  \item 3.车端面B4.车外圆G
  \item 5.切断工序 02：车削
  \item 仅需1步：车端面C（简单端面修整）
  \item 工序 03：铣削安装：1次 (回转夹具)
  \item 工位：2个 (利用夹具回转)
  \item 1.铣侧面D (工位1)
  \item 2.铣侧面E (工位2，回转后)
  \item 工件的装夹方式找正装夹法
  \item 夹具装夹法成批、大量生产首选
  \item 获得加工尺寸精度的四种方法试切法
  \item 通过“试切-测量-调整-再试切”的反复过程获得尺寸。
  \item 预先调整好刀具与工件的相对位置，加工中保持不变。
  \item 直接利用刀具的相应尺寸来保证工件被加工尺寸。
  \item 例子：钻头、铰刀加工自动获得法
  \item 完全由机床设备本身的精度来保证获得规定精度。
  \item 和进度计划。
  \item 式中 Q——产品的年产量(台/年)；
  \item n——每台产品中该零件的数量(件/台)；
  \item a％——备品率；
  \item b％——废品率。
  \item 案例：重型机械、专用设备、新品试制。
  \item 案例：机床厂、电机厂。
  \item 案例：汽车厂、轴承厂、电视机流水线。
  \item 机械加工工艺规程工艺规程的核心作用
  \item 指导生产：指挥现场生产操作的主要技术依据管理依据：用于技术准备、生产调度及成本核算建厂基础：规划车间面积、设备购置与布局的资料技术交流：工艺技术人员沟通与协作的标准介质工艺规程包括：
  \item 工艺路线拟订零、部件在生产过程中，各部门或工序的先后顺序
  \item 机械加工工艺规程的制定机械加工零件结构的工艺性
  \item 研究分析零件图及车间的具体情况，准备原始资料了解现有的生产条件
  \item 确定各个工序的加工余量、工序尺寸及公差。
  \item 确定各工序的切削用量及工时定额。
  \item 按照加工顺序标注尺寸，避免多尺寸同时保证。
  \item 从定位基准标注尺寸，避免基准不重合误差。
\end{itemize}

\subsection{第二节 定位基准的选择}
\begin{tcolorbox}[colback=yellow!6,colframe=red!55!black,title=重点,left=1mm,right=1mm]
\begin{enumerate}[leftmargin=2.2em,label=\arabic*.]
  \item 定义：设计图样上所采用的基准，是确定零件各表面间相对位置的依据。
  \item 核心原则：实际生产中应力求基准重合（设计基准与工艺基准一致），以减少加工误差基准：用来确定生产对象上几何要素间的几何关系所依据的那些点、线、面机械制造技术基础C 第八章 机械加工工艺规程的制订
  \item \textbf{基准 – 选择原则}：编制零件的加工工艺规程时，应尽量选用设计基准作
  \item \textbf{为工序基准}：加工及测量工作时，应尽量选用工序基准作为定位基
  \item \textbf{准及测量基准}：设计机器零件时，应尽量选用装配基准作为设计基准以减少加工误差以消除由于基准不重合引起的误差
  \item \textbf{一次。}：粗基准：用毛坯上未经机械加工过的表面作为定位基准或基面的称为粗基准
  \item \textbf{加工。}：精基准：用经过机械加工过的表面作为定位基准或基面的称为精基准
\end{enumerate}
\end{tcolorbox}
\noindent\textbf{次重点与补充：}
\begin{itemize}[leftmargin=2.2em,label=--]
  \item 基准设计基准
  \item 典型案例：轴类零件的中心线，是各外圆表面的设计基准，决定了圆柱面的同轴度。
  \item 工艺基准在工艺过程中采用的基准称为工艺基准
  \item 定位基准加工中确定工件位置
  \item 工序基准工序图上用来确定本工序被
  \item 加工表面加工后的尺寸、形状、位置的基准
  \item 测量基准检验时的测量依据
  \item 装配基准装配时的位置依据
  \item 定位基准的选择 – 粗基准以不加工表面为基
  \item 准保证不加工表面与加工
  \item 表面之间的位置精度，确保零件非加工部分的
  \item 外观质量。
  \item 余量最小表面优先优先选择余量最小的表
  \item 面，防止因余量不足导致该表面无法加工，造
  \item 成零件报废。
  \item 重要表面为基准确保重要表面的加工余
  \item 量均匀，从而保证其加工精度和表面质量，避
  \item 免余量分布不均。
  \item 平整光洁表面选择面积大、无浇口/冒
  \item 口/飞边的表面，保证定位准确、夹紧可靠，避
  \item 免因定位不稳产生误差。
  \item 同一方向只使用一次
  \item 粗基准精度低，重复使用会产生较大的定位误
  \item 差，导致后续加工精度超差，因此通常只使用
  \item 定位基准的选择 – 精基准基准重合原则
  \item 尽量选择设计基准或工序基准作为定位基准，从源头避
  \item 免基准不重合带来的加工误差。
  \item 基准统一原则尽可能多的表面采用同一精
  \item 基准加工，既能保证位置精度，又能简化夹具设计与制
  \item 造。
  \item 互为基准原则当两个表面相互位置精度要
  \item 求很高时，可互为基准反复加工，逐步提高加工精度。
  \item 自为基准原则对于精度要求很高且余量小
  \item 而均匀的表面，直接以加工表面自身作为定位基准进行
\end{itemize}

\subsection{第三节 机械加工工艺路线的拟订}
\begin{tcolorbox}[colback=yellow!6,colframe=red!55!black,title=重点,left=1mm,right=1mm]
\begin{enumerate}[leftmargin=2.2em,label=\arabic*.]
  \item 加工方法的选择加工经济精度
  \item 正常条件下，每种加工方法所能保证的精度等级与
  \item 表面粗糙度范围。
  \item \textbf{IT6级以上高精度}：由粗到精，逐步提高精度与表面质量，合理分配加工余量与切削负荷
  \item 核心：追求高生产率任务：使主要表面达到一定
  \item 核心：精度过渡任务：保证各主要表面符合
  \item 核心：最终精度保障任务：解决表面质量问题（
  \item 核心：IT6级以上高精度划分加工阶段的目的与价值
  \item \textbf{及时发现毛坯缺陷}：及早发现夹渣、气孔等问题，避免浪费后
  \item \textbf{合理插入热处理工序}：便于在粗精加工之间安排退火、淬火等工序，改善材料性能与加工性。
  \item \textbf{- 定义：少数工序完成复杂任务，内容高度集成}：- 优点：减少安装次数，保证精度，效率高（如加工中心）
  \item - 缺点：设备复杂昂贵，生产准备工作量大工序分散
  \item \textbf{- 定义：工序数量多，每道工序仅完成少量简单操作}：- 优点：设备工装简单，调整维修方便，适应性强- 缺点：工序流转多，生产组织管理相对复杂决策依据：根据生产批量大小、设备条件及产品精度要求灵活选择方案
  \item \textbf{先基准后其他}：先加工精基准，再以其定位加工其他表面
  \item \textbf{先粗后精}：粗加工在前快速去除余量，精加工在后保证精度
  \item \textbf{先主后次}：优先完成主要表面，次要表面穿插进行
  \item \textbf{先面后孔}：先加工平面，为孔加工提供稳定可靠的基准02 热处理与辅助工序时机退火/正火：安排在粗加工前或后调质处理：安排在粗加工之后进行淬火处理：安排在精加工之前进行
  \item \textbf{表面处理：安排在工艺过程的最后阶段}：检验工序：关键工序、阶段转换、车间转换之后
\end{enumerate}
\end{tcolorbox}
\noindent\textbf{次重点与补充：}
\begin{itemize}[leftmargin=2.2em,label=--]
  \item 典型表面的加工方法的选择
  \item 针对外圆、孔、平面等不同表面特征，以及尺寸大
  \item 小选择适配方法机械制造技术基础C 第八章 机械加工工艺规程的制订
  \item 加工阶段的划分粗加工阶段
  \item 高效去除大部分余量，接近零件形状。
  \item 追求高生产率半精加工阶段
  \item 使主要表面达到一定精度，为精加工做准备。
  \item 精度过渡精加工阶段
  \item 保证各主要表面符合图样规定的质量要求。
  \item 最终精度保障光整加工阶段
  \item 解决表面质量问题（如降低粗糙度）。
  \item 任务：高效去除大部分余量，接近零件形状。
  \item 精度，为精加工做准备。
  \item 图样规定的质量要求。
  \item 如降低粗糙度）。
  \item 续精加工时间与成本。
  \item 保证加工表面质量消除粗加工产生的变形和应力影响，在后
  \item 续精加工阶段得到修正。
  \item 合理选择机床设备粗加工使用大功率低精度机床；
  \item 精加工使用高精度机床，优化资源配置。
  \item 工序的集中与分散工序集中
  \item 加工序的安排01 加工顺序四大原则
\end{itemize}

\subsection{第四节 加工余量及工序尺寸和公差的确定}
\begin{tcolorbox}[colback=yellow!6,colframe=red!55!black,title=重点,left=1mm,right=1mm]
\begin{enumerate}[leftmargin=2.2em,label=\arabic*.]
  \item \textbf{关键影响因素解析}：表面质量：需覆盖粗糙度及缺陷层深度尺寸公差：补偿上道工序的尺寸偏差范围位置误差：如弯曲、偏斜等形位误差需余量抵
  \item \textbf{消}：安装误差：本工序定位基准与装夹带来的偏差$\Delta Z_{\text{total}}=\Delta Z_1+\Delta Z_2+\Delta Z_3+\cdots+\Delta Z_n$
  \item \textbf{公差依据：经济精度}：根据各工序加工方法的经济精度等级来确定工序尺寸
  \item \textbf{标注规范：入体原则}：外表面尺寸注负偏差，内表面尺寸注正偏差，确保加
\end{enumerate}
\end{tcolorbox}
\noindent\textbf{次重点与补充：}
\begin{itemize}[leftmargin=2.2em,label=--]
  \item 加工余量的确定：加工总余量（$\Delta Z_{\text{总}}$）为毛坯与成品设计尺寸之差，即切除金属层总厚度。
  \item 工序余量（$\Delta Z_i$）为相邻两工序尺寸之差，即单工序切除厚度。
  \item 相邻两工序尺寸之差，即单工序切除厚度。
  \item - 单边余量：平面加工- 双边余量：轴/孔、回转面
  \item ni<1
  \item ΔZi机械制造技术基础C 第八章 机械加工工艺规程的制订
  \item 工序尺寸和公差的确定推算逻辑：由后向前
  \item 从零件的设计尺寸出发，逐步向前推算出各工序尺寸及毛坯尺寸。
  \item 公差。
  \item 工余量充足。
\end{itemize}

\subsection{第五节 工艺尺寸链}
\begin{tcolorbox}[colback=yellow!6,colframe=red!55!black,title=重点,left=1mm,right=1mm]
\begin{enumerate}[leftmargin=2.2em,label=\arabic*.]
  \item 尺寸链中的每个尺寸称为尺寸链的环封闭环 (A0)：自然形成的一环在加工或装配过程最后形成，其尺寸是间接保证的，尺寸精度取决于所有组成环。
  \item \textbf{组成环：影响封闭环的全部环}：- 增环 (Ai↑)：该环增大，封闭环也增大- 减环 (Ai↓)：该环增大，封闭环减小
  \item \textbf{增减环判断法：回路箭头法}：从封闭环出发沿回路绕行，箭头方向与封闭环相反则为增环，方向相同则
  \item \textbf{设计尺寸链}：由零件设计图上的尺寸所组成的尺寸组合
  \item \textbf{空间尺寸链}：尺寸链中各环不在同一平面或彼此平行的
  \item 平面内按应用场合分： 按空间位置分：
  \item \textbf{04 封闭环公差}：ES0 ＝ ESp－ EiqEI0 ＝ EIp － ESq
  \item \textbf{L1=40:0.05}：:0.09mm= 40.07±0.02mm
  \item \textbf{EI1=ES2−ES0}：EI1=0−+0.25 mm=−0.25mm
  \item \textbf{;0.15mm}：按入体原则标注，则为 A1=39.85;0.1
  \item \textbf{0 mm。}：尺寸链的计算 – 测量基准与设计基准不重合时
  \item \textbf{其中 40;0.3}：0 mm为间接得到的，为封闭环L0
  \item \textbf{L1为减环。}：L1= 40−10mm=30mmT1= 0.3−0.16mm=0.14mm
  \item \textbf{=9.8mm}：ES0=ES1+ES3−EI2ES1=ES0+EI2−ES3=0.18+−0.09−0mm
  \item \textbf{=0.09mm}：EI0=EI1+EI3−ES2EI1=EI0+ES2−EI3=−0.18+0−−0.09 mm
  \item \textbf{=−0.09mm}：所以 A1=9.8mm±0.09mm
\end{enumerate}
\end{tcolorbox}
\noindent\textbf{次重点与补充：}
\begin{itemize}[leftmargin=2.2em,label=--]
  \item 尺寸的概念什么是尺寸链？
  \item 相互联系且按一定顺序排列的封闭尺寸组合
  \item 为减环。
  \item 机械制造技术基础C 第八章 机械加工工艺规程的制订尺寸链的分类
  \item 工艺尺寸链加工过程的各有关工艺尺寸所组成的尺寸
  \item 链称为工艺尺寸链装配尺寸链
  \item 机器装配过程中，由相互连接的尺寸形成封闭的尺寸组合称为装配尺寸链
  \item 称为设计尺寸链线性尺寸链
  \item 全部组成环均平行于封闭环平面尺寸链
  \item 组成环均位于一个或几个平行平面内，各环之间可以不平行
  \item 尺寸链的计算方法 – 极值法01 封闭环基本尺寸
  \item 封闭环的公称尺寸等于各增环公称尺寸之和减去各减环公称尺寸之和
  \item 02 环的极限尺寸03 封闭环极限偏差
  \item kp 1
  \item mkq 1
  \item A0＝ Ap－ AqAmax= A+ES
  \item Amin = A+EIT0 =ES0 - EI0
  \item 尺寸链的计算 - 定位基准与设计基准不重合时尺寸 L1是本工序保证的，尺寸 100
  \item :0.18mm 要通过前面三个尺寸间接保证，所以 L0=100
  \item :0.18mm是封闭环。
  \item L1、L2为增环， L3为减环L0=L1+L2−L3
  \item 10mm=L1+30mm−60mmL1=40mm
  \item ES0=ES1+ES2−EI30.18mm=ES1+0.04mm−−0.05mmES1=+0.09mm
  \item EI0=EI1+EI2−ES30mm=EI1+0mm−+0.05mmEI1=0.05mm
  \item 尺寸链的计算 – 定位基准与设计基准不重合时A0400
  \item :0.25mm 为封闭环，A1为减环，
  \item A280;0.150 mm 为增环
  \item A0=A2−A140mm=80mm−A1
  \item A1=40mmES0=ES2−EI1
  \item EI0=EI2−ES1ES1=EI2−EI0
  \item ES1= −0.15−0mm=−0.15mmA1=40;0.25
  \item 而 40;0.160 mm为增环，
  \item ES1=0.14mmEI1=0
  \item 所以，尺寸 L1=300:0.14mm。
  \item 尺寸链的计算 – 需继续加工的表面标注的工序尺寸计算画尺寸链
  \item A0是间接形成的封闭环，A1、A3为增环，
  \item A2为减环。
  \item 公称尺寸：
  \item A0=A1+A3−A2A1=43.4mm
  \item 上极限偏差：
  \item ES0=ES1+ES3−EI2ES1=+0.315mm
  \item 下极限偏差：
  \item EI0=EI1+EI3−ES2EI1=+0.05mm
  \item 所以 A1=43.4:0.05:0.315mm
  \item 入体尺寸原则写为 A1=43.450:0.265mm。
  \item 尺寸链的计算 – 对多尺寸保证的工艺尺寸链计算3工序钻 4工序车 5工序磨A010mm±0.18mm 为封闭环，A1、A3为增环, A2为减环。
  \item A1=A0−A3+A2= 10−9.2+9mm
  \item 机械制造技术基础C 第九章 工件在机床上的安装
\end{itemize}

\subsection{第一节 概述}
\begin{tcolorbox}[colback=yellow!6,colframe=red!55!black,title=重点,left=1mm,right=1mm]
\begin{enumerate}[leftmargin=2.2em,label=\arabic*.]
  \item 工件的安装核心概念解析
  \item \textbf{工件安装}：包含定位(确定正确位置) 与夹紧(固定位置防
  \item 位移)两个过程。
  \item \textbf{机床夹具}：机床上用于装夹工件的工艺装备，是实现精准
  \item \textbf{生产效率}：缩短辅助工时，实现快速装夹与更换，显著提升加工节拍。
  \item \textbf{劳动条件}：降低工人操作难度与劳动强度，减少对高技能水平的依赖。
  \item \textbf{加工成本}：减少废品率与试切损耗，优化工序流程，有效控制综合成本
  \item 优点：方法简单，无需专用夹具，成本低。
  \item 缺点：效率低，精度依赖工人技术水平。
  \item 应用：单件小批生产专用夹具安装
  \item 优点：精度高，效率高，减轻劳动强度。
  \item 缺点：夹具设计制造周期长，成本较高。
  \item 应用：中批及大批大量生产夹具的组成
  \item \textbf{接所有元件。}：核心三要素：定位元件、夹紧装置、夹具体这三部分是任何机床夹具必不可少的核心组成部分，缺一不可。
  \item \textbf{专用夹具}：针对特定工件工序设计。效率极高，成本较高，适合中批以
  \item \textbf{组合夹具}：由标准件拼装而成，灵活多变。可快速拆卸重组，适合新产
  \item \textbf{成组夹具}：适应成组工艺，通过调整个别元件，可加工形状相近的多种工件。
  \item \textbf{随行夹具}：自动线专用，承载工件依次通过各加工工位，实现自动化流转加工。
\end{enumerate}
\end{tcolorbox}
\noindent\textbf{次重点与补充：}
\begin{itemize}[leftmargin=2.2em,label=--]
  \item 安装的关键工具。
  \item 定位 + 夹紧 = 正确安装加工质量
  \item 直接决定尺寸精度、形状及位置精度，是加工的基础保障。
  \item 。
  \item 机械制造技术基础C 第九章 工件在机床上的安装找正安装
  \item 依据工件表面或划线痕，用划针/指示表逐个找正位置后夹紧。
  \item 靠夹具保证工件与刀具及机床的相对位置，从而直接保证加工精度。
  \item 1. 定位元件确定工件位置，如V形块、定位
  \item 销、支承钉等。
  \item 2. 夹紧装置紧固工件防止移动，如压板、螺
  \item 旋夹紧、液压缸。
  \item 3. 对刀/导向装置确定刀具位置，如对刀块、钻套
  \item 、镗套。
  \item 4. 连接元件确定夹具相对机床位置，典型如
  \item 定向键。
  \item 5. 其他元件根据需求配置，如分度装置、顶
  \item 出器等辅助件。
  \item 6. 夹具体夹具的基础骨架，用于安装和连
  \item 夹具的分类通用夹具
  \item 机床标准附件，如卡盘、平口钳。通用性强，适合单件小批生产。
  \item 上生产。
  \item 品试制。
\end{itemize}

\subsection{第二节 工件定位原理}
\begin{tcolorbox}[colback=yellow!6,colframe=red!55!black,title=重点,left=1mm,right=1mm]
\begin{enumerate}[leftmargin=2.2em,label=\arabic*.]
  \item 六点定位原理核心：六点定位原理
  \item \textbf{完全定位}：工件的6个自由度全部被限制，是最严谨的定位
  \item \textbf{方式。}：应用：矩形工件钻孔、连杆小头孔加工
  \item \textbf{不完全定位}：根据加工要求，无需限制工件的全部自由度。
  \item 应用：车削轴径（无需限制轴向移动/转动）
  \item \textbf{欠定位}：应限制的自由度没有被限制，无法保证加工精度
  \item \textbf{过定位}：多个支承点重复限制同一个或几个自由度。
\end{enumerate}
\end{tcolorbox}
\noindent\textbf{次重点与补充：}
\begin{itemize}[leftmargin=2.2em,label=--]
  \item 通过合理布置六个定位支撑点，限制工件的六个自由度，使工件在夹具中占据唯一确定的位置，是定位设计的基本准则。
  \item 工件的六个自由度物体在空间具有六个自由
  \item 度，即沿三个坐标轴的移动(分别用
  \item 表示)和绕三个坐标轴的转表示)。
  \item 定位 != 夹紧定位解决“位置确定”问题；
  \item 夹紧解决“防止受力移动”问题。
  \item 两者概念不同，缺一不可。
  \item 机械制造技术基础C 第九章 工件在机床上的安装定位的四种状态
  \item 。
  \item ?? 严重影响质量，绝对不允许出现！
  \item 判断：需具体分析，有时可提高刚度，非一概禁止。
  \item 定位方式与定位元件的选择
\end{itemize}

\subsection{第三节 定位方式与定位元件的选择}
\begin{tcolorbox}[colback=yellow!6,colframe=red!55!black,title=重点,left=1mm,right=1mm]
\begin{enumerate}[leftmargin=2.2em,label=\arabic*.]
  \item \textbf{锥套}：机械制造技术基础C 第九章 工件在机床上的安装常见典型定位方式及定位元件所限制的自由度
\end{enumerate}
\end{tcolorbox}
\noindent\textbf{次重点与补充：}
\begin{itemize}[leftmargin=2.2em,label=--]
  \item 常见典型定位方式及定位元件所限制的自由度工件定位基准面 定位元件
  \item 平面支撑钉
  \item 支撑板固定支撑与自位支撑
  \item 固定支撑与辅助支撑圆孔
  \item 定位销锥销
  \item 锥孔顶尖
  \item 锥心轰外圆柱面
  \item 支撑钉戒支撑板V形块
  \item 定位套半圆孔
  \item 短销
  \item 长机械制造技术基础C 第九章 工件在机床上的安装
\end{itemize}

\subsection{第四节 定位误差}
\begin{tcolorbox}[colback=yellow!6,colframe=red!55!black,title=重点,left=1mm,right=1mm]
\begin{enumerate}[leftmargin=2.2em,label=\arabic*.]
  \item \textbf{基准不重合误差 ($\Delta B$)}：由于工序基准和定位基准不重合而引起的误差，$\Delta B$ 的大小为从定位基准到设计基准之间在加工尺寸方向上的最大变化量在该方向上的投影。
  \item \textbf{基准位移误差 ($\Delta Y$)}：由于定位元件或工件制造不准确，导致定位基准在工序尺寸方向上发生位置偏移而产生的误差。
  \item \textbf{定位误差公式}：$\Delta D=\Delta B+\Delta Y$（矢量和）。
\end{enumerate}
\end{tcolorbox}
\noindent\textbf{次重点与补充：}
\begin{itemize}[leftmargin=2.2em,label=--]
  \item 定位误差及其组成矢量关系判定法则
  \item 同向相加变动方向对工序尺寸影响相同
  \item （如均变大），误差求和。
  \item 反向相减变动方向对工序尺寸影响相反
  \item （一增一减），误差取差。
  \item 五步分析法1. 确定工序尺寸与基准
  \item 2. 确定定位基准3. 计算基准不重合误差
  \item 4. 计算基准位移误差5. 判定方向，合成总误差
\end{itemize}

\section{第十章 机械加工精度}
\subsection{本章核心知识}
\textit{（本节暂无可提取文本）}

\subsection{第一节 概述}
\begin{tcolorbox}[colback=yellow!6,colframe=red!55!black,title=重点,left=1mm,right=1mm]
\begin{enumerate}[leftmargin=2.2em,label=\arabic*.]
  \item 机械加工精度的基本概念加工精度 vs 加工误差
  \item 精度是实际与理想参数的符合程度零件加工后实际几何参数与理想几何参数接近程度。
  \item 指加工后的实际几何参数（尺寸，形状和表面间的相互位置）对理想几何参数的偏离程度尺寸精度
  \item \textbf{02 成形法}：利用成形刀具的轮廓直接“复印”出工件形状，加工效率
  \item \textbf{03 仿形法}：刀具严格按照仿形装置（如靠模）的轮廓轨迹进行进给加
  \item \textbf{04 展成法（范成法）}：工件和刀具做特定的展成运动（啮合运动），加工复杂曲
  \item \textbf{多次安装获得法}：工件有关表面的位置精度，完全依赖于加工表面与定位基准面之间的位置关系来间接保证。
\end{enumerate}
\end{tcolorbox}
\noindent\textbf{次重点与补充：}
\begin{itemize}[leftmargin=2.2em,label=--]
  \item 加工表面本身的尺寸（如直径、长度）与
  \item 其理想尺寸的符合程度。
  \item 形状精度加工表面宏观几何
  \item 形状（如圆度、平面度）的误差
  \item 位置精度加工表面与其他表
  \item 面之间的位置关系（如平行度、垂直
  \item 度）的误差什么是经济精度？
  \item 指在正常生产条件下（采用标准设备、工艺装备及技术等级工人，不延长加工时间），某种加工方法所能稳定保证的加工精度。
  \item 获得尺寸精度的方法试切法
  \item 通过“试切-测量-调整-再试切”的反复过程来获得尺寸
  \item 精度。适用于单件小批量生产。
  \item 调整法先调整好刀具与工件的相对
  \item 位置，并在一批零件的加工过程中保持不变。适用于成
  \item 批大量生产。
  \item 定尺寸刀具法利用刀具的相应尺寸来保证
  \item 工件的加工尺寸。例如：钻孔、铰孔加工场景。
  \item 自动获得法通过自动控制系统在加工过
  \item 程中实时测量工件尺寸，并自动补偿刀具位置偏差。
  \item 获得形状精度的方法01 机床运动轨迹法
  \item 利用机床成形运动，使刀尖与工件的相对运动轨迹形成所需表面形状。
  \item 高，适用于特定曲面。
  \item 工。
  \item 面（如齿轮）。
  \item 获得位置精度的方法一次安装获得法
  \item 在工件的同一次装夹中，集中加工所有有相互位置要求的表面，无需二次定位。
  \item 位置精度极高，彻底避免了多次安装带来的基准不重合误差与定位误差。
  \item 受机床运动精度、工件装夹精度以及夹具定位精度等多重因素制约。
\end{itemize}

\subsection{第二节 工艺系统几何误差对加工精度的影响}
\begin{tcolorbox}[colback=yellow!6,colframe=red!55!black,title=重点,left=1mm,right=1mm]
\begin{enumerate}[leftmargin=2.2em,label=\arabic*.]
  \item 工艺系统-机床误差-主轰径向圆跳动
  \item 主轴回转轴线相对于平均回转轴线在径向的变动量，主要影响工件圆度、圆柱度
  \item \textbf{若机床主轴采用滚动轴承结构：}：在车床上车外圆时，滚动轴承内环外滚道的圆
  \item \textbf{度误差对主轴径向圆跳动影响较大。}：在镗床上镗孔时，轴承外环内滚道的圆度误差
  \item \textbf{3. 两导轨的平行度（扭曲）}：当前后导轨在垂直平面内 有平行度误差（扭曲误差）
  \item \textbf{工艺系统-机床误差-传动链}：传动链误差是指传动链始末两端传动元件间相对运动
  \item \textbf{的误差}：有些加工方法（如车螺纹，滚齿，插齿等），要求刀
  \item \textbf{误差来源构成}：包括定位误差、夹紧误差、安装误差以及夹具的磨损等
\end{enumerate}
\end{tcolorbox}
\noindent\textbf{次重点与补充：}
\begin{itemize}[leftmargin=2.2em,label=--]
  \item 对主轴径向圆跳动影响较大。
  \item 轴向窜动中心线沿轴线方向变动，影响端面垂直度、平面度
  \item Delta\textasciitilde{}0Delta
  \item 止推轴承端面误差对主轴轴向窜动的影响
  \item 角度摆动中心线倾斜摆动，直接影响工件圆柱度和端面的形状误差
  \item 工艺系统-机床误差-导轨1. 垂直平面内的直线度
  \item 导轨在垂直平面内有直线度误差 $\Delta z$ 时，也会使车刀在垂直面内发生位移，使工件半径产生误差 $\Delta R$，且 $\Delta R\approx(\Delta z)^2/(2R)$。
  \item 与 $\Delta z$ 相比，$\Delta R$ 属于微小量，一般可忽略不计。2. 水平面内的直线度（核心）
  \item 导轨在水平面内有直线度误差 $\Delta y$ 时，工件半径产生 $\Delta R=\Delta y$ 的误差。该误差处于敏感方向，影响极大，直接决定加工尺寸精度。
  \item 时，刀尖的运动轨迹是一条空间曲线，使工件产生圆柱度误差。
  \item 导轨间在垂直方向有平行度误差 $\Delta l_3$ 时，将使工件与刀具的正确位置在误差敏感方向产生 $\Delta y\approx(H/B)\times\Delta l_3$ 的偏移量，使工件半径产生 $\Delta R=\Delta y$ 的误差，对加工精度影响较大。
  \item 具与工件之间必须具有严格的传动比关系。
  \item 机床传动链误差是影响这类表面加工精度的主要原因之一
  \item 提高传动元件制造精度和装配精度，减少传动件数，均可减小传动链误差工艺系统-刀具 工艺系统-夹具
  \item 刀具误差定尺寸刀具（钻头/铰刀）
  \item 刀具本身的尺寸精度直接决定工件的加工尺寸精度成形刀具（成形车刀）
  \item 刀具的轮廓形状直接复制给工件，决定工件形状精度普通单刃刀具（车刀）
  \item 制造精度影响较小，但刀具磨损会显著改变切削位置夹具误差
  \item 对工件的影响直接影响工件在机床上的位置精度，进而影响形状与尺
  \item 寸精度夹具精度设计原则
  \item 工件精度（IT5\textasciitilde{}IT7），则夹具精度通常需取工件精度的1/5\textasciitilde{} 1/3
  \item 工件精度IT8，则夹具精度通常需取工件精度的 1/10\textasciitilde{} 1/5
\end{itemize}

\subsection{第三节 工艺系统受力变形对加工精度的影响}
\begin{tcolorbox}[colback=yellow!6,colframe=red!55!black,title=重点,left=1mm,right=1mm]
\begin{enumerate}[leftmargin=2.2em,label=\arabic*.]
  \item 系统产生的弹性与塑性变形加工后果
  \item 刚度定义：弹性系统抵抗变形的能力。
  \item 公式：K = F / Y (N/mm)
  \item 柔度定义：刚度的倒数，单位力产生的变形。
  \item 公式：W = 1 / K (mm/N)
  \item \textbf{Kg}：Kst、 Kjc、 Kj、 Ka、 Kg——工艺系统、机床、夹具、
  \item \textbf{误差复映}：待加工表面上有什么样的误差，加工表面上必然也有同样性质的误差，这就是切削加工中的误差复映现象
  \item \textbf{误差复映系数}：$\varepsilon$ 称为误差复映系数，其值为小于 1 的正数。
  \item 内应力重新分布对加工精度的影响内应力：
  \item - 切削加工：塑性变形与热效应影响与控制对策
  \item \textbf{影响：导致工件缓慢变形，丧失加工精度。}：对策：精密零件必须进行时效处理（自然/人工）以
  \item - 结构设计优化：采用封闭截面、增设肋板，提升零部件自身刚性。
  \item \textbf{减小工艺系统外力}：- 降低切削用量：采用较小的背吃刀量和进给量，从源头
  \item \textbf{减小切削力。}：- 刀具几何优化：合理选择刀具角度（如90°偏刀），有
\end{enumerate}
\end{tcolorbox}
\noindent\textbf{次重点与补充：}
\begin{itemize}[leftmargin=2.2em,label=--]
  \item 基本概念- 受力变形现象外力来源
  \item 切削力、夹紧力、传动力及惯性力变形本质
  \item 破坏刀具工件位置，产生误差基本概念- 刚度和柔度
  \item 数值越大，系统越稳定。
  \item 数值越小，加工精度越高。
  \item 工艺系统总变形系统总变形 =
  \item 机床变形 + 夹具变形 + 刀具变形 + 工件变形Yst=Yjc+Yj+Yd+YgYst—工艺系统的变形量；
  \item Yjc—机床的变形量；
  \item Yj—夹具的变形量；
  \item Yd—刀具的变形量；
  \item Yg—工件的变形量Kst= 1
  \item Kje+1
  \item KjKd
  \item 刀具、工件的刚度。
  \item 工艺系统刚度的测定与分析刚度测定-单向加载测定
  \item 在车床的两顶尖间，装一根刚性很好的心轴1，并在刀架上装加力器5，对正心轴1的中点位置，4为测力环，
  \item 通过螺钉8加力Fy，由千分表9的读数可知 Fy的大小(测力环需标定)，
  \item 由千分表2、千分表3、千分表7分别读出机床主轴箱、尾座和刀架的变形量。
  \item 已知刀架处的作用力为Fy，作用于前、后顶尖处的作用力为Fy /2
  \item 工艺系统受力变形对加工精度的影响-切削力切削力作用点位置变化
  \item 在加工过程中，刀具相对于工件的位置是不断变化的，也就是说，切削力的作用点位置和切削力的大小是在变化的。
  \item 同时，工艺系统在各作用点上的刚度一般是不相同的机床变形
  \item 工件变形工艺系统变形
  \item Yg=Fy3EI
  \item L−x2x2L
  \item Yst=Yjc+Yg=Fy
  \item Kdj+
  \item KctL;x
  \item Kwzx
  \item L;x2x23EIL
  \item Yjc=Fy切削力大小变化
  \item 加工余量不均或材料硬度不一致，也会影响工件的加工精度
  \item 由于背吃刀量变化引起切削力变化，工艺系统变形也产生相应变化
  \item 可进行多次进给加工来减小误差复映工艺系统受力变形对加工精度的影响-其他作用力
  \item 01 夹紧力产生的加工误差针对刚性差的工件（如薄壁套），过大的夹紧力会产生弹性形变，加工后松开夹具即回弹，造成形状误差02 惯性力产生的加工误差不平衡部件高速旋转产生的离心力方向周期性变化，导致工艺系统受力状态波动，引发工件加工误差。
  \item 03 传动力产生的加工误差如单爪拨盘驱动工件旋转时，其分力方向随旋转周期性改变，对工艺系统变形的影响机理与惯性力类似。
  \item 外部载荷去除后，残存在工件内部的应力。零件内部存在内应力时，其内部组织处于一种极不稳定的平衡状态，这种组织强烈地要求恢复到一个没有应力的稳
  \item 定状态，即使在常温下，零件也会缓慢地发生这种变化，直到内应力消失为止
  \item 内应力主要来源：
  \item - 毛坯制造：冷却不均- 工件冷校直：强制变形
  \item 消除内应力。
  \item 减少工艺系统受力变形的措施提高工艺系统刚度
  \item - 接触刚度提升：改善接触面粗糙度，施加预紧力消除配合间隙。
  \item - 工件安装优化：增加跟刀架/中心架等辅助支承，合理选择装夹方式。
  \item 效减小径向切削分力。
\end{itemize}

\subsection{第四节 工艺系统受热变形对加工精度的影响}
\begin{tcolorbox}[colback=yellow!6,colframe=red!55!black,title=重点,left=1mm,right=1mm]
\begin{enumerate}[leftmargin=2.2em,label=\arabic*.]
  \item 工艺系统的热源内部热源
  \item \textbf{切削热（核心热源）}：金属弹塑性变形与摩擦产生，占比最大。
  \item \textbf{摩擦热}：导轨、轴承、齿轮等运动副摩擦发热。
  \item \textbf{动力源发热}：电动机、液压系统运行时的能量损耗。
  \item \textbf{外部热辐射}：来自阳光直射、照明灯光及取暖设备的辐射热
  \item 工艺系统热变形对加工精度的影响机床热变形：
  \item - 主轴系统：主轴箱发热导致轴线抬高/倾斜
  \item - 床身导轨：温差引起床身弯曲变形- 后果：工件产生圆柱度、锥度等形状误差工件热变形：
  \item - 均匀受热：轴类车削，引起尺寸热胀冷缩- 不均受热：薄板磨削，单面受热导致翘曲- 后果：冷却后产生平面度或尺寸精度偏差减小工艺系统热变形的措施
\end{enumerate}
\end{tcolorbox}
\noindent\textbf{次重点与补充：}
\begin{itemize}[leftmargin=2.2em,label=--]
  \item 外部热源环境温度变化
  \item 受季节更替、昼夜温差影响，具有周期性波动特征。
  \item 量。
  \item 减少热源合理控制切削用量，优
  \item 化刀具参数，并改善运动副的摩擦特性，从源
  \item 头降低发热量。
  \item 隔离热源将电机、油箱等强热源
  \item 移出主机；对无法移出的发热部件，采用隔热
  \item 材料进行有效隔离。
  \item 加强冷却采用大流量切削液冲刷
  \item 、高压冷却或喷雾冷却技术，强制带走加工区
  \item 域产生的大量热量。
  \item 保持热平衡精加工前让机床空运转
  \item 预热，使其快速达到热稳定状态，减小加工过
  \item 程中的动态变形。
  \item 控制环境温度对于超精密加工，必
  \item 须在恒温车间进行，严格控制环境温度波
  \item 动对工件精度的影响。
\end{itemize}

\section{第十二章 装配工艺}
\subsection{本章核心知识}
\textit{（本节暂无可提取文本）}

\subsection{第一节 装配工艺的定制}
\begin{tcolorbox}[colback=yellow!6,colframe=red!55!black,title=重点,left=1mm,right=1mm]
\begin{enumerate}[leftmargin=2.2em,label=\arabic*.]
  \item 为半成品或成品的工艺过程。
  \item \textbf{装配工艺的核心任务}：通过分析零件精度与产品精度的关系，选择合理的装配方法与过程，确保达到
  \item \textbf{01 清洗作业}：装配前对零件进行擦洗、浸洗或喷洗，确保洁净度，保障质量与寿命。
  \item \textbf{02 部件连接}：分为螺纹、销等可拆卸连接，以及焊接、铆接等不可拆卸连接方式。
  \item \textbf{03 校正与调整}：找正零部件相互位置，调整运动副间隙，确保装配精度与运行顺畅。
  \item \textbf{04 配作加工}：以已加工件为基准，加工配合件或进行组合加工，保证配合精度。
  \item \textbf{05 动/静平衡}：对高速回转件进行静平衡或动平衡调试，消除离心力，防止振动。
  \item \textbf{2. 减少手工劳动量}：针对刮研等传统高手工环节，大力推进装配机械化与自动化，以此提升生产效率并缩短装配周期。
  \item \textbf{3. 降低装配成本}：通过优化资源配置，合理选用流水线或自动线，在减少设备投资的同时，最大化利用车间生产面积。
  \item 核心依据：
  \item 艺文件单件小批：装配系统图
  \item 成批生产：装配工艺过程卡大批大量：详细装配工序卡
  \item \textbf{大数互换法}：基于概率解法，允许0.27\%的不合格品。可显著放宽零件公差，降低加工难度与成本，适用于大批量生产。
  \item 核心思想：当装配精度极高（互换法不经济）时，先将组成环公差放大至经济加工精度，再通过选择合适的零件配对来满足装配要求。
  \item \textbf{直接选配法}：工人凭经验直接挑选零件进行装配。效率较低，对工人技
  \item \textbf{分组选配法}：将零件实测尺寸分组，同组内零件可互换。适用于大批量生产（
  \item \textbf{保证装配精度的方法 – 修配法}：各组成环按“经济精度”加工，装配时通过修配尺寸链中某一指定环（修配环）的尺寸，补偿累积误差，从而
  \item 核心思想：装配时通过改变产品中可调件的相对位置或选用调整件，在不切削加工的前提下达到装配精度要求。
  \item \textbf{可动调整法}：改变零件相对位置，如拧动螺钉、调整楔块，实时消除累积误差。
  \item \textbf{固定调整法}：加入不同尺寸的调整件（如垫片、衬套），按需选择补偿尺寸链误差。
  \item \textbf{误差抵消调整法}：合理选择装配方向，利用误差的方向性，使部分误差相互抵消以提高精度
\end{enumerate}
\end{tcolorbox}
\noindent\textbf{次重点与补充：}
\begin{itemize}[leftmargin=2.2em,label=--]
  \item 装配工作的内容什么是装配？
  \item 按产品规定的技术要求，将合格零件或部件进行配合和连接，使之成
  \item 部件装配将若干个零件装配成部件的过程，是
  \item 总装配前的关键环节。
  \item 总装配把零件和部件装配成最终产品的过程
  \item ，直接决定产品的最终性能。
  \item 产品规定的精度要求。
  \item 06 最终检验按技术标准对装配成品进行性能指标
  \item 检验，并进行试运转测试。
  \item 定制装配工艺的原则1. 确保产品质量
  \item 装配是产品质量的最终保障环节，同时也是发现零件加工或设计缺陷、促进质量闭环改进的关键机会。
  \item 定制装配工艺的步骤01. 研究装配图纸和
  \item 验收技术要求- 深度解析总装与部件图纸，理
  \item 清结构关系- 明确产品工作原理与性能指标
  \item - 审查验收技术标准，预判并解决关键难点
  \item 02. 确定装配生产组织形式
  \item - 固定式：工位固定，适合单件/小批生产
  \item - 移动式：流水线作业，适合大批/大量生产
  \item （注：移动式装配生产效率更高，成本更低）
  \item 03 划分装配单元与装配顺序
  \item 装配单元：分解为合件、组件等，便于流水作业。
  \item 装配顺序：遵循“先下后上、先内后外、先难后易
  \item ”原则。
  \item 04 合理选择装配方法
  \item 1. 装配精度要求2. 产品结构特点
  \item 3. 生产纲领规模05 编制装配工
  \item 保证装配精度的方法 – 互换法核心思想
  \item 核心优势：装配高效。
  \item 操作简单，生产率高生产协作
  \item 便于组织流水线生产维护便捷
  \item 利于产品维修更换完全互换法
  \item 基于极值解法，保证100\%零件互换。
  \item 当封闭环公差要求高且组成环多时，对零件加工精度要求极高，成本较高。
  \item 保证装配精度的方法 – 选择装配法
  \item 术水平依赖度高。
  \item 如活塞与销）。
  \item 使封闭环达到规定精度。
  \item 应用场景- 成批或单件小批量生产模式
  \item - 装配精度要求极高，互换法无法满足的场景
  \item 修配环选择原则- 易于修配且装卸方便
  \item - 避免选用并联尺寸链的公共环- 非表面处理（如淬火）的零件保证装配精度的方法 – 调整法
\end{itemize}

\section{第十三章 先进制造技术}
\subsection{本章核心知识}
\textit{（本节暂无可提取文本）}

\subsection{第一节 工业4.0与智能制造}
\begin{tcolorbox}[colback=yellow!6,colframe=red!55!black,title=重点,left=1mm,right=1mm]
\begin{enumerate}[leftmargin=2.2em,label=\arabic*.]
  \item \textbf{德国国家级战略背景}：2013年汉诺威工业博览会正式推出，旨在利用前沿技术重塑制造业竞争力，引领第四次工业革命浪潮。
  \item \textbf{核心：信息物理系统 (CPS)}：打通供应、制造、销售全链路，将物理设备与数字世
\end{enumerate}
\end{tcolorbox}
\noindent\textbf{次重点与补充：}
\begin{itemize}[leftmargin=2.2em,label=--]
  \item 工业4.0工业1.0
  \item 蒸汽机时代蒸汽动力机械化
  \item 工业2.0电气化时代
  \item 电机与流水线生产工业3.0
  \item 信息化时代IT技术实现自动化
  \item 智能化时代CPS驱动产业变革
  \item 界深度融合，实现生产过程的全面数据化与智慧化。
  \item 模式转变从集中式控制向分散式增强型控
  \item 制转变核心目标
  \item 建立高度灵活的个性化、数字化产品服务模式
  \item 技术核心基于CPS实现智能工厂，动态
  \item 配置柔性生产智能工厂
  \item 研究智能化生产系统及过程，实现网络化分布式生产设施的高效协作
  \item 智能生产涵盖生产物流管理、人机互动及3D技术
  \item 应用，优化生产流程效率智能物流
  \item 通过互联网、物联网深度整合物流资源，大幅提升供应链响应速度
  \item 智能制造智能装备
  \item 工业机器人、智能数控系统等硬件基础
  \item 智能系统体现数字化、网络化的核心制造系
  \item 统智能产品
  \item 实现全生命周期个性化定制与数据互联
  \item 智能服务贯穿产品全生命周期的主动式增值
  \item 服务智能工厂：落地载体
  \item 通过构建智能化生产系统与网络化生产设施，实现生产过程的全面智能化。它是将智能制造技术应用于实际生产的物理空间。
  \item 自主能力自主感知与决策
  \item 整体可视全流程数据透明
  \item 自我学习持续优化与进化
  \item 人机交互高效协作与互动
\end{itemize}

\subsection{第二节 增材制造技术}
\begin{tcolorbox}[colback=yellow!6,colframe=red!55!black,title=重点,left=1mm,right=1mm]
\begin{enumerate}[leftmargin=2.2em,label=\arabic*.]
  \item 增材制造核心原理：离散-堆积
  \item \textbf{01 三维建模}：使用CAD软件设计零件的三维数字模型
  \item \textbf{02 分层切片}：将三维模型分层，获取各层二维截面信息
  \item \textbf{03 数控加工}：生成数控代码，精确控制设备逐层加工
  \item \textbf{04 层层叠加}：连续加工薄层并自动连接，最终形成实体
  \item 材料：液态光敏树脂原理：紫外激光逐点扫描固化
  \item 特点：精度高、表面光洁，成本较高选择性激光烧结 (SLS)
  \item 材料：尼龙、金属、陶瓷粉末原理：激光扫描粉末层使其烧结连接
  \item 特点：材料广泛，无需支撑，精度稍低
  \item 材料：纸、塑料薄膜等薄片原理：激光切割轮廓 + 热压辊粘合
  \item 特点：成形快，不易变形，无需支撑熔融沉积制造 (FDM)
  \item 材料：ABS、PLA等丝状材料原理：喷头熔化丝状材料并逐层挤出
  \item 特点：成本低、结构简单，桌面级主流
  \item 材料：金属或合金粉末原理：高能激光完全熔化粉末，逐层
  \item 堆积特点：制造高性能金属零件，接近全
  \item 材料：金属熔丝原理：利用电弧熔化金属丝材进行沉
  \item 积特点：沉积效率高，适合大型金属零
\end{enumerate}
\end{tcolorbox}
\noindent\textbf{次重点与补充：}
\begin{itemize}[leftmargin=2.2em,label=--]
  \item 技术核心优势高度柔性化
  \item 数字化制造，制造成本与批量大小无关快速制造
  \item 大幅减少加工工序，显著缩短产品周期自由成形
  \item 不受专用工具限制，完美加工复杂异形零件材料广泛
  \item 支持树脂、金属、陶瓷等多种材料成形光固化成形 (SLA)
  \item 分层实体制造 (LOM)
  \item 选择性激光熔化 (SLM)
  \item 致密电弧增材制造 (WAAM)
  \item 件增材制造技术的应用领域
  \item 产品创新设计快速制作高精度模型，
  \item 辅助新产品开发与设计验证，缩短研发迭代周
  \item 期。
  \item 快速铸造技术直接制造复杂铸件的蜡
  \item 型和砂型，显著缩短模具试制周期，提升铸造
  \item 效率。
  \item 高端零部件制造制造航空航天等领域高
  \item 性能、难加工的复杂核心零部件，实现轻量化
  \item 与功能集成。
  \item 生物医学工程定制化制造植入体、个
  \item 性化假体及组织工程支架，助力精准医疗与康
  \item 复工程。
  \item 微纳电子制造制造微纳尺度的功能性
  \item 器件与封装结构，推动电子信息产业向微型化
  \item 发展。
  \item 增材制造技术的发展趋势复合材料化
  \item 从单一材料向多功能纳米、纤维增强复合材料
  \item 方向发展，拓展应用边界。
  \item 多工艺融合融合多种增材工艺，打
  \item 破单一技术局限，实现高质量产品的一体化精
  \item 密成形。
  \item 极限尺寸突破双向突破尺寸限制：向
  \item 米级大型化结构件与微米/纳米级精密器件同步
  \item 过程智能化融入传感器与微处理器
  \item ，实现成形过程数据的实时记录、分析与闭环
  \item 反馈。
  \item 模式融合创新与传统减材制造互补融
  \item 合，结合云制造与工业互联网，重塑柔性生产
  \item 新模式。
\end{itemize}

\section{第十四章 绿色制造技术}
\subsection{本章核心知识}
\textit{（本节暂无可提取文本）}

\subsection{第一节 绿色制造的基本概念}
\begin{tcolorbox}[colback=yellow!6,colframe=red!55!black,title=重点,left=1mm,right=1mm]
\begin{enumerate}[leftmargin=2.2em,label=\arabic*.]
  \item 绿色制造的概念与内涵核心定义与目标
  \item \textbf{定义：一种系统地考虑环境影响和资源效率的}：现代制造模式，追求全生命周期的可持续性。
  \item \textbf{目标：产品全周期环境负面影响最小化，资源}：利用效率最大化，兼顾经济效益与生态效益。
  \item \textbf{双目标导向}：1. 减少污染，保护环境；2. 优化资源利用，提升效率，实现双赢。
\end{enumerate}
\end{tcolorbox}
\noindent\textbf{次重点与补充：}
\begin{itemize}[leftmargin=2.2em,label=--]
  \item 多学科交叉融合制造技术、环境科学与资源利
  \item 用等多领域知识，形成系统性解决方案。
  \item 双过程驱动覆盖产品全生命周期过程与物
  \item 流转化过程，确保流程闭环高效。
  \item 核心技术支撑以绿色设计为源头，清洁生产
  \item 为过程，绿色再制造为延伸。
\end{itemize}

\subsection{第二节 绿色设计}
\begin{tcolorbox}[colback=yellow!6,colframe=red!55!black,title=重点,left=1mm,right=1mm]
\begin{enumerate}[leftmargin=2.2em,label=\arabic*.]
  \item 绿色设计概念和内涵核心概念解析
  \item 定义：系统地考虑环境影响，并将其集成到产品最初设计过程中的技术和方法。
  \item 核心：在满足功能、质量和成本的同时，优化设计因素，使产品全生命周期对环境影响最小
\end{enumerate}
\end{tcolorbox}
\noindent\textbf{次重点与补充：}
\begin{itemize}[leftmargin=2.2em,label=--]
  \item 。
  \item Reduce 减量化从源头减少环境污染和
  \item 能源消耗，精简结构，提高利用率。
  \item Reuse 再利用直接利用产品或零部件
  \item ，通过维修、翻新等方式延长其使用寿命。
  \item Recycle 再循环产品废弃后拆解分类，
  \item 回收原材料进行再生循环使用，变废为宝。
  \item 绿色设计策略优选环保材料
  \item 选用清洁、可再生、低能耗及高可回收性的原材料。
  \item 减少材料用量优化结构设计，最大限度减小
  \item 产品重量与体积。
  \item 优化生产技术采用绿色工艺，精简加工步骤
  \item ，降低生产能耗。
  \item 优化销售系统推行无包装/简包装设计，采
  \item 用高能效的物流运输方案。
  \item 降低使用影响优化产品能效，鼓励使用清
  \item 洁能源，减少使用阶段排放延长产品寿命
  \item 提升可靠性与耐用性，设计为易于维护、升级的形态。
  \item 优化回收系统采用模块化与易拆卸设计，
  \item 最大化材料回收与重复利用率。
\end{itemize}

\subsection{第二节 清洁生产}
\begin{tcolorbox}[colback=yellow!6,colframe=red!55!black,title=重点,left=1mm,right=1mm]
\begin{enumerate}[leftmargin=2.2em,label=\arabic*.]
  \item 清洁生产核心定义与理念
  \item 【实质】物料和能耗最少的生产规划。将废物减量化、资源化和无害化，从生产过程中消除污染。
  \item 【核心】坚持“预防为主”，从源头削减污染，彻底改变“末端治理”的被动模式。
  \item \textbf{2. 减少排放相容}：降低工业废弃物与污染物排放，推动工业生产与生态环境
  \item 干式加工技术-核心定义：加工全程摒弃切削液，通过优化刀具与工
  \item \textbf{-核心效益：}：节约成本（切削液占制造成本16\%），减少污染，改
  \item 近净成形技术-核心定义：零件成形后仅需少量加工或无需加工即可
  \item \textbf{直接装配使用。}：大幅减少原材料浪费，降低后续加工能耗，显著提升
\end{enumerate}
\end{tcolorbox}
\noindent\textbf{次重点与补充：}
\begin{itemize}[leftmargin=2.2em,label=--]
  \item 核心实施目标1. 合理利用资源
  \item 综合利用各类资源，落实节能、降耗、节水措施，最大化提升资源利用率。
  \item 和谐共存。
  \item 艺实现高效加工。
  \item 善车间环境。
  \item -关键技术：干式车削、干式滚切、干式磨削等。
  \item 生产效率。
  \item -关键技术：精密铸造、精密锻压、精密焊接等。
\end{itemize}

\subsection{第五节 绿色制造技术的意义与前景}
\begin{tcolorbox}[colback=yellow!6,colframe=red!55!black,title=重点,left=1mm,right=1mm]
\begin{enumerate}[leftmargin=2.2em,label=\arabic*.]
  \item -核心价值：重大意义推动制造科学发展
  \item \textbf{实现可持续发展}：落实人类可持续战略，实现经济与环境协调发展。
  \item \textbf{显著经济效益}：降低生产与环境成本，提升企业核心市场竞争力。
  \item \textbf{突破贸易壁垒}：提升产品绿色属性，从容应对国际市场准入挑战。
  \item -未来展望：应用前景国家战略高度
  \item \textbf{国际标准驱动}：ISO 14000环境管理体系标准的普及，强力推动绿色
  \item \textbf{新兴市场机遇}：催生绿色设计、清洁生产等全新产业链与市场增长点
\end{enumerate}
\end{tcolorbox}
\noindent\textbf{次重点与补充：}
\begin{itemize}[leftmargin=2.2em,label=--]
  \item 变革现代制造思想，完善可持续制造理论体系。
  \item 《中国制造2025》将“绿色制造工程”列为五大重点实施工程之一。
  \item 技术的广泛落地。
  \item 。
\end{itemize}

\end{spacing}

\end{document}
